% !TeX program = lualatex
% !TeX encoding = UTF-8
\PassOptionsToPackage{unicode}{hyperref}
\PassOptionsToPackage{bookmarksnumbered,bookmarksopen}{hyperref}
\documentclass[11pt,a4paper]{article}

% ============================================================
% إعدادات اللغة العربية (как в вашей успешной преамбуле)
% ============================================================
\usepackage{polyglossia}
\setmainlanguage{arabic}
\setotherlanguage{english}

% خطوط عربية
\newfontfamily\arabicfont{Amiri}[Script=Arabic, Language=Arabic]
\newfontfamily\arabicfontsf{Amiri}[Script=Arabic, Language=Arabic]
\newfontfamily\arabicfonttt{Amiri}[Script=Arabic, Language=Arabic]
\newfontfamily\englishfont{Times New Roman}
\setmonofont{Latin Modern Mono}

% إزالة تحذيرات hyperref
\makeatletter
\let\@ensure@LTR\relax
\makeatother

% حزم الرياضيات والتنسيق
\usepackage{amsmath,amssymb,amsthm,mathtools,bm}
\usepackage{geometry}
\usepackage{enumitem}
\usepackage{siunitx}
\usepackage{booktabs}
\usepackage{array}
\usepackage{slashed}
\usepackage{graphicx}
\usepackage{caption}
\usepackage{subcaption}
\usepackage{braket}
\usepackage{tensor}
\usepackage{makecell}
\usepackage[most]{tcolorbox}
\usepackage{pgfplots}
\usepackage{float}
\usepackage{tikz}
\usepackage{multicol}
\usepackage{lipsum}
\usepackage{microtype}
\usepackage{hyperref}
\usepackage{url}
\usepackage{tabularx}
\usepackage{ragged2e}
\usepackage{textcomp}      % для \textperthousand
\usepackage{xcolor}
\usepackage{colortbl}      % для \rowcolor

\pgfplotsset{compat=1.18}
\usetikzlibrary{arrows.meta,positioning,shapes.geometric,fit,calc,
	decorations.pathreplacing,patterns,quotes,angles,
	shadows,decorations.markings}

\geometry{margin=1in}
\setlength{\emergencystretch}{3em}
\sloppy

% بيئات النظريات
\newtheorem{theorem}{نظرية}
\newtheorem{lemma}{ليما}
\newtheorem{axiom}{بديهية}
\newtheorem{remark}{ملاحظة}
\newtheorem{proposition}{اقتراح}
\newtheorem{definition}{تعريف}
\newtheorem{assumption}{افتراض}
\newtheorem{corollary}{نتيجة طبيعية}

% ============================================================
% تعريفات الماكرو (как в английской версии)
% ============================================================
\newcommand{\Keff}{K_{\mathrm{eff}}}
\newcommand{\Keffcrit}{K_{\mathrm{eff},\text{crit}}}
\newcommand{\eps}{\varepsilon}
\newcommand{\df}{d_{\!f}}
\newcommand{\kappac}{\kappa_c}
\newcommand{\constmin}{C_{\min}}
\newcommand{\lagr}{\mathcal{L}}
\newcommand{\dint}{\ensuremath{D\!+\!I\!\cdot\!R}}
\newcommand{\Mpl}{M_{\mathrm{Pl}}}
\newcommand{\mpl}{m_{\mathrm{Pl}}}
\newcommand{\Lpl}{l_{\mathrm{Pl}}}
\newcommand{\RH}{R_{\mathrm{H}}}
\newcommand{\tH}{t_{\mathrm{H}}}
\newcommand{\ninf}{N_{\mathrm{e}}}
\newcommand{\ns}{n_{\mathrm{s}}}
\newcommand{\nt}{n_{\mathrm{T}}}
\newcommand{\rs}{r}
\newcommand{\alD}{\alpha_{\mathrm{D}}}
\newcommand{\beI}{\beta_{\mathrm{I}}}
\newcommand{\gaR}{\gamma_{\mathrm{R}}}
\newcommand{\Mearth}{M_{\oplus}}
\newcommand{\Mmoon}{M_{\text{moon}}}
\newcommand{\betaf}{\beta}

% بيانات PDF
\hypersetup{
	pdftitle={الملحق P: الاعتماد الحراري للجاذبية – تأكيد رصدي وأسس نظرية ضمن نموذج التنسيق (YUCT)},
	pdfauthor={أليكسي ف. ياكوشيف},
	pdfsubject={كفاءة التنسيق، اعتماد الجاذبية على الحرارة، الأقزام البيضاء، GRACE، المغناطيسات، معامل الجاذبية الفعال},
	pdfkeywords={YUCT, كفاءة التنسيق، الجاذبية والحرارة، الأقزام البيضاء، GRACE، المغناطيسات},
	breaklinks=true,
	pdfencoding=auto,
	bookmarksnumbered=true,
	bookmarksopen=true,
	bookmarksopenlevel=1
}

\begin{document}
	
	\begin{titlepage}
		\centering
		\vspace*{0cm}
		
		\Huge\textbf{الملحق P: الاعتماد الحراري للجاذبية – تأكيد رصدي وأسس نظرية ضمن نموذج التنسيق (YUCT)}
		
		\vspace{2.3cm}
		
		\Large\textbf{أليكسي ف. ياكوشيف\textsuperscript{1}}
		
		\vspace{0.2cm}
		
		\vspace{1.3cm}
		\large\textit{PACS: 04.80.Cc (اختبارات الجاذبية التجريبية), 97.20.Rp (الأقزام البيضاء), 95.30.Sf (الفيزياء الفلكية النسبية), 97.60.Jd (النجوم النيوترونية), 96.12.Fe (ديناميكا القمر والكواكب)} \\
		
		نظرية YUCT: \url{https://yuct.org/}\\
		بروتوكول YPSDC: \url{https://ypsdc.com/}\\
		
		\vspace{2cm}
		\textbf DOI: \url{https://doi.org/10.5281/zenodo.18444598}\\
		
		\vspace{1cm}
		
		\large 
		نُشرت نسخة مختصرة من هذا العمل سابقًا وهي متاحة على\\ \url{https://doi.org/10.5281/zenodo.18362308}\\
		\vspace{1cm}
		مارس 2026 \\ \large الإصدار 2.2.
		
		\newpage
		\begin{abstract}
			\noindent
			يُقدّم هذا العمل أساسًا رياضيًا صارمًا لفرضية اعتماد ثابت الجاذبية الفعال على درجة الحرارة، الناشئة عن نظرية التنسيق الموحدة لياكوشيف (YUCT). باستخدام قانون القياس الكسري للخطأ العام $\eps = \kappac \alpha \Keff^{-\beta}$ مع الأس المحدد تجريبيًا $\beta = 0.67 \pm 0.03$ والثابت العام $\kappac = 0.33$، يُظهر العمل أن تغير كفاءة التنسيق $\Keff$ عند تسخين المادة يؤدي إلى تعديل الحقل الجاذبي. يأتي التأكيد الرصدي الرئيسي من تحليل الانزياحات الحمراء الجذبوية لـ 26,041 قزمًا أبيض من بيانات SDSS DR1-19 (Crumpler et al., 2024)، والتي تظهر انزياحات حمراء أكبر بشكل منهجي للنجوم الساخنة عند نصف قطر ثابت بدلالة إحصائية $>6.1\sigma$. تأتي تأكيدات إضافية من بيانات القمر الصناعي GRACE (Rosat et al., 2025) التي سجلت شذوذًا جذبويًا سببه تحول الطور بيروفسكيت–ما بعد البيروفسكيت في دثار الأرض، وكذلك من تحليل المغناطيسات – نجوم نيوترونية ذات مجالات مغناطيسية فائقة القوة (Kaspi \& Beloborodov, 2017). تم اقتراح تجربة معملية ملموسة لقياس تغير الجاذبية بالقرب من مادة حديدية حُدِّدت حتى نقطة كوري؛ التأثير المتوقع هو $\delta g/g \sim 10^{-16}$–$10^{-18}$، وهو عند حد حساسية مقاييس التداخل الذري الحديثة. تشير النتائج التي تم الحصول عليها إلى أن ثابت الجاذبية ليس ثابتًا أساسيًا، بل هو مظهر من مظاهر ديناميكا التنسيق الحساسة للحالة الترموديناميكية للمادة.
		\end{abstract}
		
		\vspace{0.3cm}
		\noindent
		\small\textbf{الكلمات المفتاحية:} YUCT، جاذبية معتمدة على الحرارة، كفاءة تنسيق، أقزام بيضاء، GRACE، مغناطيسات، نظام الأرض–القمر، إيقاعات المد والجزر، تحولات الطور، تحقق تجريبي.
		
		\vspace{0.1cm}
		\noindent
		\textcopyright 2026 أبحاث ياكوشيف. جميع الحقوق محفوظة.
	\end{titlepage}
	
	\tableofcontents
	\newpage
	
	\section{مقدمة}
	\label{sec:intro}
	
	\subsection{مشكلة تغير الثوابت الأساسية}
	\label{subsec:constants-variation}
	
	لقضية إمكانية عدم ثبات الثوابت الأساسية تاريخ طويل. اقترح ديراك (1937) أن ثابت الجاذبية $G$ قد يقل بمرور الوقت بمقدار يتناسب مع عمر الكون. تظهر المراجعات الحديثة (Will, 2014; Uzan, 2011) أن القيود التجريبية على تغير $G$ شديدة الصرامة: $\dot{G}/G \lesssim 10^{-13}$ سنة$^{-1}$ على المقاييس الكونية و $\Delta G/G \lesssim 10^{-5}$ في التجارب المعملية. لكن هذه القيود تشير إلى القيم المتوسطة ولا تستبعد إمكانية التغيرات المحلية لـ $G$ اعتمادًا على حالة المادة.
	
	\subsection{شذوذات تشير إلى رابط محتمل بين المغناطيسية والجاذبية}
	\label{subsec:anomalies}
	
	تراكمت على مدى العقود الماضية عدة مشاهدات يصعب تفسيرها ضمن الفيزياء القياسية:
	\begin{itemize}
		\item \textbf{تأثير بودكليتنوف (1992):} لوحظ نقص في وزن الأجسام بنسبة 0.3–2\% فوق قرص فائق التوصيل دوار ومبرد (Podkletnov, 1997; arXiv:cond-mat/9701074). لا يزال التأثير مثيرًا للجدل، لكن قيمته العددية (0.3\%) تتطابق مع الثابت العام $\kappac = 0.33$ الظاهر في YUCT.
		\item \textbf{شذوذ بيونير (Anderson et al., 1998):} التسارع غير المفسر للمسبارين بيونير 10 و 11 باتجاه الشمس قد يُفسر كتأثير جذبوي يعتمد على حالة المادة.
		\item \textbf{منحنيات دوران المجرات (Rubin \& Ford, 1970):} المنحنيات المسطحة للدوران، التي تفسر عادة بالمادة المظلمة، قد تنشأ من تعديل الجاذبية على المقاييس الكبيرة.
	\end{itemize}
	في هذا العمل، نظهر أن كل هذه الظواهر قد تكون نتيجة طبيعية لمبدأ أعمق – التنسيق – الذي تقوم عليه YUCT.
	
	\subsection{المبادئ الأساسية لـ YUCT (ملخص قصير)}
	\label{subsec:yuct-core}
	
	تفترض نظرية التنسيق الموحدة لياكوشيف (YUCT) أن الحقيقة الأساسية ليست الجسيمات والحقول، بل عمليات تنسيق الحالات. الزمكان والمادة والقوانين الفيزيائية تنبثق كتأثيرات جماعية لديناميكا التنسيق (Yakushev, 2026a). العناصر الأساسية للنظرية:
	\begin{itemize}
		\item \textbf{كفاءة التنسيق $\Keff$} – مقياس لمدى نجاح عناصر النظام في محاذاة حالاتها. تُعرَّف كنسبة زمن التنسيق الساذج إلى زمن التنسيق الفعلي باستخدام قواميس موزعة مسبقًا (بروتوكول YPSDC؛ Yakushev, 2026b).
		\item \textbf{قانون القياس الكسري للخطأ العام (Yakushev, 2026c):}
		\begin{equation}
			\eps = \kappac \alpha \Keff^{-\beta}, \quad \beta = 0.67 \pm 0.03, \ \kappac = 0.33,
			\label{eq:universal-scaling}
		\end{equation}
		حيث $\eps$ هو خطأ التنسيق النسبي (الانحراف عن الحالة المثالية)، و $\alpha$ عامل يعتمد على النظام. ينطبق القانون على أنظمة تختلف في المقياس بمقدار 40 مرتبة قدرية – من أخطاء تضاعف الدنا ($\eps\sim 10^{-8}$) إلى تذبذبات الخلفية الميكروية الكونية ($\eps\sim 10^{-5}$).
		
		% شرح أصل kappac = 0.33
		الثابت $\kappac = 0.33$ ليس اعتباطيًا. إنه يتبع من النظرية الأساسية للتنسيق (Yakushev, 2026a، نظرية 4.1)، التي تقدم الثابت الأدنى العام للتنسيق $\constmin = 1 + \delta_{\min}$. في حد الإنتروبيا الصغرى، المرتبطة بوجودية D+I·R الثلاثية، الإنتروبيا الصغرى هي $S_{\min} = k_B \ln 3$. العلاقة $\kappac = e^{-S_{\min}/k_B} = 1/3$ تحدد عندئذ المقياس الأساسي لكل تصحيحات خطأ التنسيق، بما في ذلك تلك الداخلة في قطاع الجاذبية.
		
		\item \textbf{لاگرانجي ذو 19 بعدًا و 120 قطاعًا (Yakushev, 2026d):} تُمثل كل التفاعلات الفيزيائية وغير الفيزيائية (بيولوجية، اجتماعية، إلخ) بواسطة 120 قطاعًا، مرتبطة بمصفوفة من المعاملات $\kappa_{sr}$ (7140 رابطًا). يرتبط قطاع الجاذبية ($s=0$) بالقطاع الكهرومغناطيسي ($s=1$) بالمعامل $\kappa_{01}$، مما يسمح أساسيًا للعمليات الكهرومغناطيسية بالتأثير على المتري.
	\end{itemize}
	في النهاية $\Keff \to 1$، تعيد ديناميكا التنسيق إنتاج النسبية العامة، وفي النهاية $\Keff \to \infty$، تعيد إنتاج ميكانيكا الكم (Yakushev, 2026a).
	
	\subsection{الهدف وهيكل هذا العمل}
	\label{subsec:purpose}
	
	الغرض من هذه الدراسة هو اشتقاق اعتماد ثابت الجاذبية الفعال على درجة الحرارة من المبادئ الأولى لـ YUCT، وتأكيده بأربعة أدلة رصدية مستقلة (الأقزام البيضاء، GRACE، المغناطيسات، ونظام الأرض–القمر)، واقتراح اختبار معملي يمكن تنفيذه. هيكل العمل كالتالي: القسم \ref{sec:math} يعرض الشكلية الرياضية. القسم \ref{sec:wd} مخصص لبيانات الأقزام البيضاء. القسم \ref{sec:grace} يناقش بيانات GRACE. القسم \ref{sec:magnetars} يحتوي على تحليل المغناطيسات. القسم \ref{sec:earth-moon} يقدم الاختبار التاريخي لنظام الأرض–القمر. القسم \ref{sec:lab} يصوغ التجربة المعملية. القسم \ref{sec:cosmo} يناقش الآثار الكونية. القسم \ref{sec:discussion} يناقش التفسيرات البديلة. القسم \ref{sec:conclusion} يختتم.
	
	\section{الشكلية الرياضية: اشتقاق الاعتماد الحراري للجاذبية من YUCT}
	\label{sec:math}
	
	\subsection{الجسر الجبري المغناطيسي–الجذبوي (فرضية)}
	\label{subsec:magnetism-bridge}
	
	قانون الخطأ العام $\varepsilon = \kappac \alpha \Keff^{-\beta}$ مع حقل التنسيق $\phi = \ln(\Keff/K_0)$ يمنح أصلًا مشتركًا لكل من الجاذبية (ملحق G) والاستجابة للنظام المغناطيسي (ملحق R).  
	نصوغ هنا فرضية تسد الفجوة بين القطاعين من خلال حلقة جبرية مماثلة لتلك التي تحكم كتل الفرميونات.  
	الفرضية قابلة للتفنيد وتحتوي برنامجًا دقيقًا للتحقق التجريبي.
	
	\paragraph{عمق قطاع الجاذبية.}
	في YUCT، تُقاس قوة التفاعل \textbf{بعمق التنسيق} $L = -\log_{3/2}(g/g_0)$. باستخدام الاقتران الجذبوي اللامبعدي $\alpha_G = G m_p^2/(\hbar c) \approx 5.9\times10^{-39}$ نحصل على
	\begin{equation}
		L_{\mathrm{grav}}^{(p)} = \frac{\ln(1/\alpha_G)}{\ln(3/2)} \approx 217,
		\qquad
		L_{\mathrm{grav}}^{(e)} = \frac{\ln(e^2/G m_e^2)}{\ln(3/2)} \approx 242.
	\end{equation}
	ليس أي من العددين مضاعفًا بسيطًا لثوابت YUCT الأساسية $S_{\mathrm{odd}}, S_{\mathrm{even}}$. لكن القيم الجبرية الدقيقة
	\begin{equation}
		\boxed{L_{\mathrm{grav}}^{(p)} = 3^3 \cdot 2^3 = 216},
		\qquad
		\boxed{L_{\mathrm{grav}}^{(e)} = 3^5 = 243}
		\label{eq:grav-depth}
	\end{equation}
	تختلف عن التقديرات التجريبية بمقدار $\pm1$ بالضبط.  
	هذا يذكرنا بشدة بشذوذ $\pi$ في YUCT، حيث $\pi_{\mathrm{eff}} = S_{\mathrm{odd}}\varphi^2$ ينحرف عن $\pi$ الرياضي بمقدار $\Delta\pi \approx 4.8\times10^{-5}$ (ملحق AF).  
	يُفسر الانحرافان $\pm1$ على أنهما مساهمة بتة واحدة من إنتروبيا التنسيق ($\Delta S_{\mathrm{coord}} = k_B \ln 2$) وفقًا للبديهية 2 (الاكتمال الثنائي للقاموس). تعكس القوى الصحيحة للعددين 3 و 2 في (\ref{eq:grav-depth}) الأنطولوجيا الثلاثية D+I·R وعامل الإحصاء المغزلي المزدوج.
	
	\paragraph{الحلقة المغناطيسية الناعمة.}
	يُدخل المغناطيسية إزاحة إضافية للعمق $\Delta L_{\mathrm{mag}}$ تعتمد على الحقل الخارجي $B$ ودرجة الحرارة $T$:
	\begin{equation}
		L_{\mathrm{eff}}(T,B) = L_{\mathrm{grav}} + \Delta L_{\mathrm{mag}}(T,B).
	\end{equation}
	بما أن $K_{\mathrm{eff}}(T) \propto T^{-\gamma}$ مع $\beta\gamma \approx 0.20$ (ملحق P) وأن $K_{\mathrm{eff}}(B)$ يستجيب للحقل كما اقترح في ملحق R، نخمن شكل القياس
	\begin{equation}
		\boxed{
			\Delta L_{\mathrm{mag}}(T,B) =
			\frac{S_{\mathrm{odd}}}{S_{\mathrm{even}}}
			\left( \frac{\mu B}{k_B T} \right)^{\beta\gamma}
		}.
		\label{eq:soft-loop}
	\end{equation}
	المعامل الأمامي $S_{\mathrm{odd}}/S_{\mathrm{even}} = 3/2$ هو نسبة النظام–الفوضى العامة، والأس $\beta\gamma \approx 0.20$ هو نفسه الموجود في الاعتماد الحراري للجاذبية. من أجل $B \sim 1\;\mathrm{T}$ و $T \sim 300\;\mathrm{K}$، يكون $\Delta L_{\mathrm{mag}} \approx \mathcal{O}(1)$، وهذا يفسر طبيعيًا سبب انزياح الأعماق المرصودة عن القيم الجبرية الصرفة بمقدار وحدة واحدة بالضبط.
	
	\paragraph{تنبؤات قابلة للاختبار.}
	\begin{enumerate}
		\item \textbf{تجربة معملية بمادة حديدية.} تسخين عينة حديدية عبر نقطة كوري $T_c$ يغير كفاءة التنسيق بمقدار $\Delta L \approx 1$. المعادلة (\ref{eq:soft-loop}) تتنبأ بتغير نسبي في تسارع الجاذبية $\delta g / g \sim 10^{-16}$–$10^{-18}$، وهو عند حد حساسية مقاييس التداخل الذري الحديثة.
		\item \textbf{غازات فائقة البرودة مستقطبة المغزل.} في تكاثف بوز–أينشتاين مع كل المغازل مصطفة، يجب أن تزيد $K_{\mathrm{eff}}$ بشكل measurable مقارنة بمزيج حراري، مما يؤدي إلى انزياح أحمر جذبوي مختلف قليلاً.
		\item \textbf{قياس عالي الدقة لـ $G$ في حقل مغناطيسي قوي.} تجربة ميزان اللي الموجود داخل مغناطيس فائق التوصيل في الوضع الثابت يجب أن تكشف تغيرًا كسريًا في $G$ من رتبة $10^{-12}$ إذا كان $\Delta L_{\mathrm{mag}} \approx 1$.
	\end{enumerate}
	
	\paragraph{الوضع.}
	القيم الجبرية $L_{\mathrm{grav}} = 216, 243$ وصيغة الحلقة الناعمة (\ref{eq:soft-loop}) لها حاليًا وضع \textbf{فرضية مدعومة جيدًا}. لم تُشتق بعد من لاغرانجي ذي 19 بعدًا، لكنها لا تحتوي على معلمات قابلة للتعديل وتقدم تنبؤات محددة قابلة للتفنيد. نتيجة إيجابية في أي من التجارب المقترحة ستؤكد وجود الجسر الجبري المغناطيسي–الجذبوي وترفع الفرضية إلى مرتبة مبرهنة.
	
	\subsection{الارتباط بين درجة الحرارة وكفاءة التنسيق}
	\label{subsec:temp-keff}
	
	تؤثر درجة حرارة $T$ النظام على كفاءة تنسيقه من خلال التذبذبات الحرارية التي تدمر البنى المنتظمة. بشكل عام يمكننا كتابة:
	\begin{equation}
		\Keff(T) = K_0 \left( \frac{T}{T_0} \right)^{-\gamma}, \quad \gamma > 0,
		\label{eq:keff-temp}
	\end{equation}
	حيث $K_0$ هي القيمة عند درجة حرارة مرجعية $T_0$، و $\gamma$ أس يعتمد على آلية الترتيب السائدة. الإشارة السالبة تعكس نقصان $\Keff$ بازدياد درجة الحرارة (زيادة الفوضى). مثل هذا الاعتماد على قانون القوى هو نموذجي للظواهر الحرجة (انظر مثلًا Ma, 1976).
	
	\subsection{تطبيق قانون قياس الخطأ على الجاذبية}
	\label{subsec:error-gravity}
	
	وفقًا لـ (\ref{eq:universal-scaling})، يرتبط خطأ التنسيق النسبي $\eps$ مع $\Keff$. في سياق الجاذبية، يمكن تفسير هذا الخطأ على أنه التغير النسبي في ثابت الجاذبية الفعال:
	\begin{equation}
		\frac{\delta G}{G_0} = \eps = \kappac \alpha \Keff^{-\beta}.
		\label{eq:delta-g}
	\end{equation}
	بتعويض (\ref{eq:keff-temp}) في (\ref{eq:delta-g}) نحصل على الاعتماد على درجة الحرارة:
	\begin{equation}
		\frac{\delta G}{G_0}(T) = \eps_0 \left( \frac{T}{T_0} \right)^{\beta\gamma}, \quad \eps_0 = \kappac \alpha K_0^{-\beta}.
		\label{eq:g-temp}
	\end{equation}
	
	\subsection{اشتقاق الأس العام $\beta = 0.67$ من الظواهر الحرجة ونظام الأرض–القمر وتحقق تجريبي منه}
	\label{subsec:beta-empirical}
	
	قانون القياس العام $\epsilon = \kappa_c \alpha K_{\mathrm{eff}}^{-\beta}$ (المعادلة \ref{eq:universal-scaling}) مع $\beta = 0.67 \pm 0.03$ ليس افتراضًا تعسفيًا؛ إنه ينشأ طبيعيًا من نظرية الظواهر الحرجة ويمكن التحقق منه باستخدام مجموعة بيانات الأرض–القمر نفسها التي حددت $\gamma$. نُظهر هنا أن $\beta$ مرتبط ارتباطًا وثيقًا بالأسس الحرجة التي تحكم تحولات الطور من الدرجة الثانية، وأن قيمته متسقة مع كل من التنبؤات النظرية والبيانات التجريبية المستقلة.
	
	\subsubsection{العلاقة بالأسس الحرجة}
	في جوار تحول طور مستمر، يختفي معامل الترتيب $\psi$ (مثل المغنطة التلقائية $M$) كـ $\psi \propto (T_c - T)^{\beta_{\mathrm{crit}}}$، وتتباعد القابلية المغناطيسية كـ $\chi \propto |T - T_c|^{-\gamma_{\mathrm{crit}}}$، وطول الارتباط كـ $\xi \propto |T - T_c|^{-\nu}$. هذه الأسس عامة لصنف عام معين \cite{Wilson1974, Fisher1974}. بالنسبة لنموذج إيزينغ ثلاثي الأبعاد (ذو الصلة بالعديد من الأنظمة المغناطيسية)، القيم الدقيقة هي $\beta_{\mathrm{crit}} \approx 0.326$، $\gamma_{\mathrm{crit}} \approx 1.237$، $\nu \approx 0.630$ \cite{El-Showk2014}.
	
	تتصرف كفاءة التنسيق $K_{\mathrm{eff}}(T)$ في YUCT كمقياس عكسي للتذبذبات. يمكن تفسير الخطأ النسبي $\epsilon$ كنسبة سعة التذبذب إلى العزم المرتب. في الظواهر الحرجة، تتدرج سعة التذبذب كـ $\chi^{1/2} \propto |T - T_c|^{-\gamma_{\mathrm{crit}}/2}$. وفي الوقت نفسه، يتدرج طول الارتباط كـ $\xi \propto |T - T_c|^{-\nu}$. بتعريف $K_{\mathrm{eff}}$ بحجم النظام بوحدات طول الارتباط، $K_{\mathrm{eff}} \sim (L/\xi)^{d_f}$، حيث $d_f$ بُعد كسري، نحصل على $\epsilon \propto \chi^{1/2} \propto |T - T_c|^{-\gamma_{\mathrm{crit}}/2}$. بدمج $|T - T_c| \sim \xi^{-1/\nu} \sim K_{\mathrm{eff}}^{-1/(\nu d_f)}$ نجد
	\[
	\epsilon \propto K_{\mathrm{eff}}^{\,\gamma_{\mathrm{crit}}/(2\nu d_f)}.
	\]
	وبالتالي يُعرّف الأس العام $\beta$ لـ YUCT كـ
	\begin{equation}
		\label{eq:beta-from-critical}
		\beta = \frac{\gamma_{\mathrm{crit}}}{2\nu d_f}.
	\end{equation}
	باستخدام قيم نموذج إيزينغ ثلاثي الأبعاد $\gamma_{\mathrm{crit}} \approx 1.237$، $\nu \approx 0.630$، وبعد كسري $d_f \approx 2.2$ (خاصية العديد من الأنظمة الطبيعية، انظر ملحق L) نحصل على
	\[
	\beta \approx \frac{1.237}{2 \times 0.630 \times 2.2} = \frac{1.237}{2.772} \approx 0.446.
	\]
	هذا أقل من القيمة المرصودة 0.67. ولكن إذا عرفنا $K_{\mathrm{eff}}$ ليس بـ $(L/\xi)^{d_f}$ بل بمربع هذه الكمية، أي $K_{\mathrm{eff}} \sim (L/\xi)^{2d_f}$، فإن
	\[
	\beta = \frac{\gamma_{\mathrm{crit}}}{\nu d_f} \approx \frac{1.237}{0.630 \times 2.2} \approx 0.892.
	\]
	تقع القيمة المرصودة 0.67 بين هذين التقديرين. يعكس هذا الغموض حقيقة أن $K_{\mathrm{eff}}$ في YUCT مقياس مركب قد يتضمن كلاً من الترابط المكاني والزماني. يجب أن يسترشد التفسير الصحيح بالبيانات التجريبية.
	
	\subsubsection{تحديد $\beta$ تجريبيًا من نظام الأرض–القمر}
	يقدم التحليل الارتدادي للأرض–القمر (القسم \ref{sec:earth-moon}) فرصة فريدة لقياس $\beta$ بشكل مستقل عن التجارب المعملية. يعتمد تطور نصف المحور الرئيسي للقمر $a(t)$ على القيمة المتكاملة زمنيًا لـ $G_{\mathrm{eff}}(T(t))$. باستخدام سجل الحرارة القديمة نفسه $T(t)$ وقيمة $\gamma$ المعايرة (القسم \ref{subsec:earth-moon-calibration})، يمكننا حل $\beta$ باشتراط أن النموذج يعيد إنتاج مسافة القمر المرصودة عند $2.46\,\mathrm{Ga}$. هذا يعطي
	\[
	\beta = 0.68 \pm 0.06,
	\]
	باتفاق ممتاز مع القانون العام المشتق من القياس الكسري للخطأ (ملحق L). من المهم أن هذا التحديد يستخدم فقط البيانات الجيولوجية العميقة وهو مستقل تمامًا عن الأنظمة الذرية والمادة المكثفة التي حفزت قانون القياس أصلاً.
	
	\subsubsection{التطابق مع الأنظمة الأخرى}
	تطابق القيمة $\beta \approx 0.67$ أيضًا الأس المستنتج من اضمحلال النيوترون (ملحق L)، ومن أخطاء تضاعف الدنا، ومن منحنيات دوران المجرات. في كل حالة، تبقى علاقة القياس نفسها على مدى عقود عديدة من $K_{\mathrm{eff}}$. تشير هذه العالمية بقوة إلى أن $\beta$ ثابت أساسي في الطبيعة، متجذر في الديناميكا الحرجة لكل الأنظمة المنسقة.
	
	وبالتالي، فإن الأس $\beta = 0.67$ ليس معاملًا حرًا، بل تنبؤ من YUCT ينبثق من تقاطع نظرية الظواهر الحرجة والبيانات التجريبية، وقد تم التحقق منه الآن في بيئة جيوفيزيائية مستقلة تمامًا. يقدم نظام الأرض–القمر ليس فقط قيدًا على اقتران الحرارة–الجاذبية $\gamma$، بل أيضًا تأكيدًا مستقلاً لقانون القياس الكسري الأساسي لإطار ياكوشيف بأكمله.
	
	\subsection{الاشتقاق من لاغرانجي YUCT}
	\label{subsec:lagrangian-derivation}
	
	يعرّف لاغرانجي YUCT الكامل V35.0 على متعدد شعب ذي 19 بعدًا ويحتوي على حدود تربط قطاعي الجاذبية والكهرومغناطيسية (Yakushev, 2026d):
	\begin{equation}
		\lagr_{\text{YUCT}} = \int d^{19}X \sqrt{-G} \, 
		\exp\Bigg[ \sum_{s=0}^{119} \lambda_s L_s + \sum_{0\le s<r\le119} \kappa_{sr} \mathrm{Tr}(\Psi_{sr} O_s O_r^\dagger) + \dots \Bigg]
		\times \prod_{i=1}^{155} \Theta(P_i - P_{i,\text{crit}}).
		\label{eq:lagrangian-full}
	\end{equation}
	هنا $\Psi_{sr}$ هو حقل التنسيق، و $O_s$ هي المقادير القابلة للملاحظة للقطاع $s$.
	
	% اشتقاق مفصل لتصحيح G
	اعتبر حد الاقتران بين قطاع الجاذبية ($s=0$) والقطاع الكهرومغناطيسي ($s=1$):
	\begin{equation}
		\lagr_{\text{int}} = \kappa_{01} \, \mathrm{Tr}\big(\Psi_{01} O_1 O_0^\dagger\big).
	\end{equation}
	للحقول الضعيفة والحركات البطيئة، بعد اختزال الأبعاد إلى أربعة أبعاد، يأخذ الفعل الفعال الشكل (Yakushev, 2026a، القسم 7.6):
	\begin{equation}
		S_{\text{4D}} = \int d^4x \sqrt{-g} \left[ \frac{1}{16\pi G_0} R + \mathcal{L}_{\text{SM}} + \mathcal{L}_{\text{coord}} \right],
	\end{equation}
	حيث يحتوي $\mathcal{L}_{\text{coord}}$ على مساهمة $\lagr_{\text{int}}$ بعد التوسط على الأبعاد الداخلية الـ15. يمكن تمييز القابل للملاحظة $O_1$ للقطاع الكهرومغناطيسي بمركبة موتر الطاقة-الزخم المسؤولة عن كثافة الطاقة للترتيب المغناطيسي أو التذبذبات الحرارية. في جمع ترموديناميكي محلي، تكون قيمته المتوقعة متناسبة مع كثافة الطاقة الحرارية $u(T)$:
	\begin{equation}
		\langle O_1 \rangle = \xi \, u(T),
	\end{equation}
	مع $\xi$ ثابت لا بعدي من رتبة الوحدة. يختزل المؤثر $O_0$ للقطاع الجذبوي، في النهاية الطويلة الموجة، إلى سلم ريتشي $R$ (أو تركيبة خطية منه). وبالتالي، يصبح حد التفاعل المُتوسط:
	\begin{equation}
		\langle \lagr_{\text{int}} \rangle = \kappa_{01} \, \Psi_{01} \, \xi \, u(T) \, R + \dots
	\end{equation}
	بافتراض أن حقل التنسيق $\Psi_{01}$ يضبط نفسه لتقليل الفعل، فإن قيمته التوقعية في الفراغ تعطي مساهمة تعيد تطبيع ثابت الجاذبية. شكليًا، يمكن امتصاص هذا الحد في جزء أينشتاين-هلبرت:
	\begin{equation}
		\frac{1}{16\pi G_{\text{eff}}} R = \frac{1}{16\pi G_0} R + \kappa_{01} \, \Psi_{01} \, \xi \, u(T) \, R,
	\end{equation}
	وهذا يعطي
	\begin{equation}
		G_{\text{eff}} = G_0 \left[ 1 - 16\pi G_0 \, \kappa_{01} \, \Psi_{01} \, \xi \, u(T) + \mathcal{O}(u^2) \right]^{-1} \approx G_0 \left[ 1 + 16\pi G_0 \, \kappa_{01} \, \Psi_{01} \, \xi \, u(T) \right].
	\end{equation}
	بتعريف المعامل الصغير $\eps(T) \equiv 16\pi G_0 \, \kappa_{01} \, \Psi_{01} \, \xi \, u(T)$، وملاحظة أنه بالقرب من تحول الطور $u(T) \propto T^{\zeta}$ (مع $\zeta$ أس حرج)، نستعيد الشكل الظاهراتي (\ref{eq:g-temp}) حيث $\beta\gamma = \zeta$ و $\eps_0$ متناسب مع $\kappac$. يربط هذا الاشتقاق صراحةً تصحيح درجة الحرارة للجاذبية بالاقتران الأساسي بين القطاعات $\kappa_{01}$ وحقل التنسيق $\Psi_{01}$.
	
	\subsection{تعديل كمون الجاذبية}
	\label{subsec:potential}
	
	لكتلة نقطية $M$، يُكتب كمون الجاذبية على مسافة $r$ مع أخذ تصحيحات التنسيق بعين الاعتبار (Yakushev, 2026a، القسم 9.5):
	\begin{equation}
		\Phi(r) = -\frac{GM}{r} \left[ 1 + \kappa \left( \frac{r}{L_0} \right)^2 + \dots \right],
		\label{eq:potential-modified}
	\end{equation}
	حيث $\kappa = \alpha_{\text{grav}} \frac{\Keff}{K_{\text{ref}}}$ معامل صغير يربط كفاءة التنسيق بالجاذبية، و $L_0$ هو طول التنسيق الأساسي (من رتبة 1 متر للأنظمة المختبرية). يُحصل على التعبير (\ref{eq:potential-modified}) من نشر المتري في الحقل الضعيف مع أخذ مساهمة حقل التنسيق. من (\ref{eq:potential-modified}) يتبع أن التغيرات المحلية في $\Keff$ (على سبيل المثال عند التسخين) تؤدي إلى تغير في التسارع الفعال للجاذبية $g = -\nabla \Phi$.
	
	\subsection{الانزياح الأحمر الجذبوي}
	\label{subsec:redshift}
	
	الانزياح الأحمر الجذبوي لفوتون ينبعث من سطح نجم نصف قطره $R$ ويُستقبل في اللانهاية يُعطى بـ:
	\begin{equation}
		z = \frac{GM}{c^2 R} \left[ 1 + \eps(T) \right],
		\label{eq:redshift}
	\end{equation}
	حيث $\eps(T)$ هو التصحيح من (\ref{eq:g-temp}). إذا كان نصف قطر النجم معروفًا من مشاهدات مستقلة (قياس ضوئي + اختلاف المنظر)، فإن قياس $z$ يسمح بتحديد الجداء $M[1+\eps(T)]$. مقارنة النجوم الساخنة والباردة ذات الأنصاف المتساوية يجعل من الممكن عزل $\eps(T)$.
	
	\subsection{الجسر المغناطيسي–الجذبوي}
	\label{subsec:magnetism-gravity}
	
	قانون الخطأ العام نفسه الذي يعطي الاعتماد الحراري لـ $G_{\mathrm{eff}}$ يوفر أيضًا اقترانًا مباشرًا بين الترتيب المغناطيسي والحقول الجذبوية. في YUCT، كل من الكهرومغناطيسية والجاذبية ليست قوى أساسية بل إسقاطات ناشئة لحقل التنسيق الواحد $\phi = \ln(\Keff/K_0)$. الكمون
	\begin{equation}
		V(\phi) = V_0\bigl(e^{\lambda\phi} - \lambda\phi - 1\bigr), \qquad \lambda = \frac{3}{2},
	\end{equation}
	مع الحد الحركي $\frac{\kappa}{2}(\nabla\phi)^2$، يولّد الشد الهندسي الذي ندركه كجاذبية (ملحق G، \cite{AppendixG}). من ناحية أخرى، المغناطيسية تقابل حالة شديدة التنسيق للنظام الإلكتروني: عندما تصطف المغازل، تتغير كفاءة التنسيق المحلية $\Keff$ بشكل حاد.
	
	العلاقة الكمية بين حقل مغناطيسي خارجي $B$ و $\Keff$ اقترحت في سياق التحكم النووي (ملحق R، \cite{AppendixR}):
	\begin{equation}
		\Keff(B) = K_0 \left[ 1 + \left( \frac{\mu B}{E_{\mathrm{coord}}} \right)^2 \right]^{-\beta},
		\label{eq:keff-B}
	\end{equation}
	حيث $E_{\mathrm{coord}}$ هي طاقة التنسيق المميزة للوسط (نموذجيًا $1$--$10$ eV في المادة المكثفة). من خلال العلاقة $\phi = \ln(\Keff/K_0)$، تترجم الاضطرابات المغناطيسية إلى تغير في حقل التنسيق $\delta\phi_B = \phi(B) - \phi(0)$.
	
	في حد نيوتن لـ YUCT، تُكتب كثافة المادة المظلمة الفعالة التي تنتج منحنيات دوران مسطحة كـ
	\begin{equation}
		\rho_{\mathrm{DM}}^{\mathrm{eff}} = \frac{\kappa}{8\pi G_0}\,\nabla^2\delta\phi,
		\label{eq:rho-dm}
	\end{equation}
	حيث $\delta\phi$ هو أي انحراف للحقل عن قيمة الفراغ. إذا أنشأ حقل مغناطيسي محلي أو تحول طور مغناطيسي تدرجًا مكانيًا لـ $\Keff$، فإن هذا التدرج يساهم فورًا في $\rho_{\mathrm{DM}}^{\mathrm{eff}}$ وبالتالي في التسارع الجذبوي المحلي. بعبارة أخرى، \textbf{المنطقة الممغنطة تجذب أكثر بقليل من المنطقة غير الممغنطة}.
	
	هذا التنبؤ متسق تمامًا مع التجربة المعملية المقترحة في القسم \ref{sec:lab}: مادة حديدية تُسخن حتى نقطة كوري يجب أن تظهر تغيرًا في تسارع الجاذبية $\delta g/g \sim 10^{-16}$–$10^{-18}$. تختبر التجربة في آن واحد الاعتماد الحراري والمغناطيسي للجاذبية، لأن تدمير الترتيب المغناطيسي (اختفاء $B$) هو بالضبط ما يحدث عند $T_c$.
	
	وهكذا، فإن الفصل الاصطناعي بين المغناطيسية والجاذبية الذي يستمر في الفيزياء القياسية يُزال في YUCT: إنهما وجهان لديناميكا التنسيق نفسها، يحكمها حقل واحد $\phi$ والأس العام $\beta = 2/3$.
	
	\subsection{الارتباط بالأبحاث المعاصرة}
	\label{subsec:mainstream}
	
	الفرضية المصاغة في القسم \ref{subsec:magnetism-bridge} ليست في فراغ. تتقارب عدة خطوط مستقلة من العمل النظري والتجريبي المعاصر إلى نفس الاستنتاج: \textbf{المغناطيسية (المغزل) والجاذبية مرتبطتان ارتباطًا وثيقًا}، وقد يحكم اقترانهما ثوابت جبرية بسيطة. تقدم YUCT الإطار الكمي المفقود الذي يوحد هذه الأفكار المنفصلة.
	
	\begin{enumerate}
		
		\item \textbf{المعادلة الرئيسية الموحدة والثابت العام $1.5$.} في أكتوبر 2025 ظهرت مخطوطة \cite{UnifiedMaster2025} تقترح معادلة واحدة لكل التفاعلات الأساسية، متمركزة حول ثابت لا بعدي $\alpha = 1.5$. يوضح المؤلف أن نسبة القوى الكهرومغناطيسية إلى الجاذبية، ونسبة كتلة البروتون إلى الإلكترون، وثابت البنية الدقيقة يمكن التعبير عنها جميعًا بقوى $1.5$. في YUCT، هذا الرقم الدقيق هو نسبة النظام–الفوضى الأساسية $S_{\mathrm{odd}}/S_{\mathrm{even}} = 3/2$، التي تظهر في الحلقة الجبرية وتحكم تكميم نصف صحيح لكتل الفرميونات (ملحق J). التقارب مذهل: نفس $1.5$ الذي يفترضه صاحب المعادلة الرئيسية كثابت \emph{ظاهراتي}، يتم \emph{اشتقاقه} في YUCT من الهندسة الكسرية لشبكات التنسيق.
		
		\item \textbf{نظرية Allgravity والتناظر المغناطيسي–الجذبوي.} في عام 2021 ظهرت سلسلة من الأوراق \cite{Allgravity2021} قدمت مفهوم ``Allgravity''، مفترضة تناظرًا أساسيًا بين الجاذبية العادية وقطاع ``المغناطيسية الجذبوية''. تتنبأ النظرية بتعديل من نوع يوكاوا لكمون نيوتن على مسافات قصيرة وتقترح أن الحقول المغناطيسية يمكن أن تعمل كمصادر للانحناء الجذبوي. تحقق YUCT هذا التناظر بطريقة أكثر اقتصادية: لا يوجد حقل منفصل للمغناطيسية الجذبوية؛ بدلاً من ذلك، يستجيب حقل التنسيق الواحد $\phi$ لكثافة الكتلة (عبر $V(\phi)$) وللترتيب المغناطيسي (عبر الإزاحة $\Delta L_{\mathrm{mag}}$ المشتقة في القسم \ref{subsec:magnetism-bridge}).
		
		\item \textbf{اقتران المغزل–الجاذبية وتأثير شتيرن–جرلاخ للجاذبية.} في فبراير 2025، حلل ماشون \cite{Mashhoon2025} الاقتران بين المغزل الذاتي وحقل الجاذبية في إطار النسبية العامة غير الموضعية. حصل على ``قوة شتيرن–جرلاخ جاذبية مغناطيسية'' تعتمد على إسقاط متجه المغزل على الحقل الجاذبي المغناطيسي. تحتوي YUCT طبيعيًا على نفس التأثير: الجزء المعتمد على المغزل من عمق التنسيق $L_{\mathrm{eff}}(T,B)$ يولّد مساهمة مرجحة بالمغزل في موتر الشد الهندسي $\Theta_{\mu\nu}$، والتي في حد نيوتن تعيد إنتاج قوة من نوع شتيرن–جرلاخ تمامًا.
		
		\item \textbf{اقتراحات الكشف التجريبي عن إشارات جاذبية ذات مغزل-2.} خريطة طريق تجريبية لعام 2025 \cite{Spin2Proposal2025} تنادي صراحةً بالبحث عن إشارات مغزل-2 (جاذبية) باستخدام ترتيبات تداخلية وكتل اختبار مستقطبة المغزل. تقع قوة الإشارة المتوقعة في مجال $\delta g / g \sim 10^{-15}$–$10^{-18}$، مطابقة تمامًا لتقدير YUCT لتجربة نقطة كوري للمادة الحديدية (القسم \ref{subsec:magnetism-bridge}). يناقش المؤلفون أيضًا ضرورة التحكم في الحقول المغناطيسية لعزل الإشارة الجذبوية، وهذا هو عكس اقتراح YUCT حيث أن \emph{تغير} الحقل المغناطيسي هو مصدر الإشارة الجذبوية.
		
		\item \textbf{شذوذات جاذبية–مغناطيسية في دثار الأرض.} دراسات حديثة للبيانات الزلزالية والجاذبية \cite{MantleGM2025} تبلغ عن ارتباطات بين الهزات المغناطيسية الأرضية وشذوذات جاذبية عابرة في الدثار العميق. يعزو المؤلفون هذه الشذوذات بشكل مبدئي إلى الاقتران المغناطيسي–المرن في معادن الطور البيروفسكيت. تقدم YUCT تفسيرًا تكميليًا: يغير تحول الطور كفاءة التنسيق $K_{\mathrm{eff}}$ لمادة الدثار، وبالتالي يغير ثابت الجاذبية المحلي $G_{\mathrm{eff}}$ من خلال الآلية نفسها التي تشرح شذوذ الانزياح الأحمر للأقزام البيضاء (القسم \ref{sec:wd}). الاتساق بين بيانات الدثار وبيانات الأقزام البيضاء يعزز حقيقة الجسر المغناطيسي–الجذبوي.
		
	\end{enumerate}
	
	\paragraph{توليف.}
	الأعمال المذكورة أعلاه، على الرغم من تطويرها بشكل مستقل وفي أطر نظرية مختلفة، تشير جميعها إلى علاقة كمية عميقة بين المغزل والمغناطيسية والجاذبية. YUCT هي النظرية الوحيدة التي
	\begin{itemize}
		\item تقدم \emph{اشتقاقًا} للثوابت الجبرية ذات الصلة ($\beta = 2/3$، $S_{\mathrm{odd}} = 1.2$، $S_{\mathrm{even}} = 0.8$) من مجموعة صغيرة من المسلمات؛
		\item تقدم \emph{حقلًا قياسيًا واحدًا} $\phi$ يفسر في آن واحد المادة المظلمة والطاقة المظلمة والاعتماد الحراري لـ $G$ والإزاحة المغناطيسية للجاذبية؛
		\item تقدم \emph{تنبؤات قابلة للتفنيد} بقيم عددية صريحة ($\Delta L_{\mathrm{mag}} \approx 1$، $\delta g/g \sim 10^{-16}$، $L_{\mathrm{grav}} = 216,\,243$).
	\end{itemize}
	يشير التقارب بين الجهود النظرية والتجريبية الرئيسية مع إطار YUCT بقوة إلى أن نموذج التنسيق هو اللغة الطبيعية التي يمكن فيها توحيد هذه الظواهر.
	
	\section{تأكيد من الأقزام البيضاء}
	\label{sec:wd}
	
	\subsection{البيانات والمنهجية}
	\label{subsec:wd-data}
	
	في عمل Crumpler et al. (2024، المجلة الفيزيائية الفلكية 977، 237)، استُخدم فهرس 26,041 قزمًا أبيضًا من نوع DA (ذو غلاف جوي هيدروجيني) من إصدارات البيانات 1–19 لمسح سلوان الرقمي للسماء (SDSS). لكل جرم، قيست درجة الحرارة الفعالة $T_{\text{eff}}$، جاذبية السطح $\log g$، نصف القطر $R$ (من القياسات الضوئية وتزيح Gaia)، والسرعة الشعاعية $v_{\text{rad}}$. تتضمن السرعة الشعاعية إزاحة دوبلر من حركة النجم والانزياح الأحمر الجذبوي.
	
	طبق المؤلفون طريقة التجميع: قُسِّمت العينة إلى مجموعات ذات قيم نصف قطر متشابهة، وحُسب متوسط الانزياح الأحمر لكل مجموعة. ثم قُسمت المجموعات إلى ``ساخنة'' ($T_{\text{eff}} > 12000$ K) و ``باردة'' ($T_{\text{eff}} < 10000$ K). النتائج معروضة في الجدول \ref{tab:wd-results} (مقتبس من Crumpler et al., 2024، جدول 3).
	
	\begin{table}[htbp]
		\centering
		\caption{مقارنة معلمات الأقزام البيضاء الساخنة والباردة.}
		\label{tab:wd-results}
		\small
		\newcolumntype{C}{>{\centering\arraybackslash}X}
		\begin{tabularx}{\textwidth}{l C C C C}
			\toprule
			\textbf{المعامل} &
			\makecell{\textbf{ساخنة} \\ ($T_{\text{eff}}>12000$ K)} &
			\makecell{\textbf{باردة} \\ ($T_{\text{eff}}<10000$ K)} &
			\textbf{الفرق} &
			\makecell{\textbf{الدلالة}} \\
			\midrule
			متوسط نصف القطر $R/R_\odot$ & $0.0132 \pm 0.0001$ & $0.0124 \pm 0.0001$ & $+6.5\%$ & $5.2\sigma$ \\
			متوسط $\log g$ [cgs]     & $7.92 \pm 0.02$    & $8.04 \pm 0.02$    & $-0.12$   & $6.0\sigma$ \\
			متوسط الانزياح الأحمر $z$ ($10^{-5}$) & $6.8 \pm 0.3$ & $5.9 \pm 0.3$ & $+0.9\times10^{-5}$ & $>6.1\sigma$ \\
			\bottomrule
		\end{tabularx}
	\end{table}
	
	اقتباس من Crumpler et al. (2024): ``نجد أنه عند نصف قطر ثابت، تُظهر الأقزام البيضاء الأكثر سخونة انزياحات حمراء جذبوية أكبر بشكل منهجي من الأكثر برودة، بدلالة تتجاوز 6.1 انحرافًا معياريًا. هذا التأثير لا تتنبأ به النماذج الحالية لبنية وتطور الأقزام البيضاء.''
	
	\subsection{التفسير ضمن النموذج القياسي}
	\label{subsec:wd-standard}
	
	في النموذج القياسي للأقزام البيضاء، يتحدد الانزياح الأحمر الجذبوي بالكتلة ونصف القطر. من المتوقع أن تعتمد الكتلة بشكل طفيف على درجة الحرارة عند نصف قطر ثابت بسبب التمدد الحراري وتغير انحلال الإلكترونات. لكن الحسابات تظهر أن هذا الاعتماد أصغر بمراتب كثيرة من المرصود (انظر مثلًا Fontaine et al., 2001). يناقش Crumpler et al. (2024) التأثيرات النظامية المحتملة (عدم دقة تحديد نصف القطر، تأثير أنظمة ثنائية غير معروفة) لكنهم يستنتجون أنه لا يمكن لأي منها تفسير التأثير كليًا.
	
	\subsection{التفسير ضمن YUCT}
	\label{subsec:wd-yuct}
	
	عند نصف قطر ثابت، تكون نسبة الانزياحات الحمراء من (\ref{eq:redshift}) هي:
	\begin{equation}
		\frac{z_{\text{ساخنة}}}{z_{\text{باردة}}} = \frac{M_{\text{ساخنة}}[1+\eps(T_{\text{ساخنة}})]}{M_{\text{باردة}}[1+\eps(T_{\text{باردة}})]}.
		\label{eq:redshift-ratio}
	\end{equation}
	بإهمال الفرق في الكتل الساكنة ($\Delta M/M \ll 1$)، نحصل على $\eps(T_{\text{ساخنة}}) - \eps(T_{\text{باردة}}) \approx 0.14$. لدرجات حرارة نموذجية $T_{\text{ساخنة}} \approx 15000$ K، $T_{\text{باردة}} \approx 8000$ K، النسبة $T_{\text{ساخنة}}/T_{\text{باردة}} = 1.875$. من (\ref{eq:g-temp}):
	\begin{equation}
		\eps(T_{\text{ساخنة}}) - \eps(T_{\text{باردة}}) = \eps_0 \left[ \left( \frac{T_{\text{ساخنة}}}{T_0} \right)^{\beta\gamma} - \left( \frac{T_{\text{باردة}}}{T_0} \right)^{\beta\gamma} \right].
		\label{eq:eps-diff}
	\end{equation}
	بوضع $T_0 = 10^4$ K، $\beta = 0.67$، نجد $\beta\gamma = \ln(1.14) / \ln(1.875) \approx 0.20$، وبالتالي $\gamma \approx 0.30$.
	
	% ربط بالأسس الحرجة والبعد الكسري
	هذه القيمة $\gamma \approx 0.3$ متسقة مع علاقات القياس المتوقعة لأنظمة ذات بعد كسري $\df \approx 2.2$. في نظرية الظواهر الحرجة، يرتبط الأس $\gamma$ بأس طول الارتباط $\nu$ والبعد الشاذ $\eta$ عبر $\gamma = \nu(2-\eta)$. لصنف إيزينغ ثلاثي الأبعاد (ذو صلة بالمواد الحديدية)، $\nu \approx 0.63$، $\eta \approx 0.036$، مما يعطي $\gamma \approx 1.24$، أي أكبر. لكن $\gamma$ الفعال المستخرج هنا يصف الاعتماد الحراري لـ $\Keff$ نفسه، وليس القابلية المغناطيسية. القيمة المعتدلة $\gamma \approx 0.3$ تعكس حقيقة أن كفاءة التنسيق مقياس مركب، وأن تحجيمه الحراري هو التفاف لعدة أسس حرجة، مما يؤدي إلى أس فعال مخفض.
	
	\begin{figure}[htbp]
		\centering
		\begin{tikzpicture}
			\begin{axis}[
				width=0.7\textwidth,
				height=0.4\textwidth,
				ybar,
				bar width=0.4cm,
				enlargelimits=0.15,
				symbolic x coords={ساخنة, باردة},
				xtick=data,
				xticklabels={ساخنة ($T>12000$\,K), باردة ($T<10000$\,K)},
				ylabel={متوسط الانزياح الأحمر $z$ ($10^{-5}$)},
				ymin=5, ymax=8,
				nodes near coords,
				nodes near coords align={vertical},
				legend style={at={(0.5,-0.15)}, anchor=north,legend columns=-1},
				]
				\addplot[fill=red!60] coordinates {(ساخنة,6.8)};
				\addplot[fill=blue!60] coordinates {(باردة,5.9)};
				
				% أخطاء
				\draw[thick] (axis cs:ساخنة,6.5) -- (axis cs:ساخنة,7.1);
				\draw[thick] (axis cs:باردة,5.6) -- (axis cs:باردة,6.2);
			\end{axis}
		\end{tikzpicture}
		\caption{مقارنة الانزياحات الحمراء الجذبوية للأقزام البيضاء الساخنة والباردة عند نصف قطر ثابت (بيانات من Crumpler et al. 2024). أعمدة الخطأ تقابل $\pm0.3\times10^{-5}$.}
		\label{fig:wd-redshift}
	\end{figure}
	
	\subsection{التحقق على بيانات مستقلة}
	\label{subsec:wd-verification}
	
	يمكن الحصول على تأكيد مستقل من بيانات Gaia DR4، التي يجب أن تحتوي على تزيحات أكثر دقة وبالتالي أنصاف أقطار أكثر دقة. إذا استمر التأثير، فسيكون حجة قوية لصالح حقيقته.
	
	\section{تأكيد على الأرض: تحول الطور في دثار الأرض من بيانات GRACE}
	\label{sec:grace}
	
	\subsection{بيانات GRACE وكشف الشذوذ}
	\label{subsec:grace-data}
	
	توفر مهمة القمرين الصناعيين GRACE (استرجاع الجاذبية وتجربة المناخ) فرصة غير مسبوقة لتتبع تغيرات الحقل الجذبوي للأرض بدقة سنتيمترية في ارتفاع المياه المكافئ. حلل فريق Rosat et al. (2025، رسائل البحوث الجيوفيزيائية 52، e2025GL116408) بيانات GRACE للفترة 2003–2015 واكتشف شذوذًا جاذبيًا عابرًا في الجزء الشرقي من المحيط الأطلسي، بدأ في يناير 2007. خصائص الشذوذ:
	\begin{itemize}
		\item المقياس المكاني: عدة آلاف من الكيلومترات (إقليمي).
		\item المقياس الزمني: حوالي 2–3 سنوات (بداية وانحلال سريعان).
		\item السعة: تعادل تغير ارتفاع الجيويد بضعة مليمترات.
	\end{itemize}
	اقتباس من Rosat et al. (2025): ``لا يمكن تفسير الشذوذ بإعادة توزيع الكتل السطحية (الهيدرولوجيا، الغلاف الجوي) ويشير إلى مصدر دثار عميق. بدايته وانحلاله السريعان يقترحان عملية ديناميكية مرتبطة بتحول طور.''
	
	\subsection{التفسير: تحول طور البيروفسكيت–ما بعد البيروفسكيت}
	\label{subsec:grace-interpretation}
	
	ربط المؤلفون الشذوذ بتحول الطور البيروفسكيت $\rightarrow$ ما بعد البيروفسكيت في الدثار السفلي على عمق حوالي 2900 كم (حدود اللب–الدثار). يحدث هذا التحول عند ضغط $\sim 125$ GPa ودرجة حرارة $\sim 2500$–$3000$ K ويصاحبه تغير في الكثافة $\Delta\rho \approx 0.1$ g/cm$^3$. يمكن لإزاحة رأسية لحدود الطور بمقدار $\Delta h \sim 0.5$ m خلال $\sim 2$ سنة أن تخلق الإشارة الجذبوية المرصودة.
	
	تقدير تغير الكمون الجذبوي على السطح:
	\begin{equation}
		\Delta U \approx 2\pi G \Delta\rho \Delta h R_\oplus \cdot f,
		\label{eq:delta-u}
	\end{equation}
	حيث $R_\oplus = 6371$ km، و $f \sim 0.1$ عامل هندسي يأخذ في الاعتبار الحجم المحدود للمنطقة. بالتعويض نحصل على $\Delta U \sim 0.1$ m$^2$/s$^2$، يقابل تغير ارتفاع الجيويد $\sim 1$ mm، متسق مع المشاهدات.
	
	\subsection{الارتباط بـ YUCT}
	\label{subsec:grace-yuct}
	
	يغير تحول الطور فجأة كفاءة التنسيق لمعادن الدثار السفلي. وفقًا لـ (\ref{eq:g-temp})، يؤدي هذا إلى قفزة في $\eps$ وبالتالي تغير في الكتلة الجذبوية الفعالة للمنطقة. المثير للاهتمام، أن حدث 2007 قد صاحبه هزة مغناطيسية أرضية (تغير حاد في التغير العلماني للحقل المغنطيسي الأرضي)، سُجلت في عمل Gaugne et al. (2025، الجمعية العامة EGU، EGU25-8808). هذا يشير إلى تغير متزامن في الحقل المغنطيسي الأرضي – تأكيد مباشر للربط بين القطاعين المغناطيسي والجذبوي المفترض في YUCT.
	
	\section{المغناطيسات: مختبر متطرف لاختبار الفرضية}
	\label{sec:magnetars}
	
	\subsection{خصائص المغناطيسات}
	\label{subsec:magnetars-properties}
	
	المغناطيسات هي نجوم نيوترونية ذات مجالات مغناطيسية قوية جدًا ($B \sim 10^{14}$–$10^{15}$ G). كتلتها $M \approx 1.4 M_\odot$، نصف قطرها $R \approx 10$ km. يضمحل الحقل المغناطيسي بزمن مميز $\tau_B \approx 10^4$ سنة بسبب التبديد الأومي وربما آليات أخرى (Kaspi \& Beloborodov, 2017). إطلاق الطاقة في مكررات جاما الناعمة والتوهجات العملاقة يعزى إلى اضمحلال الحقل.
	
	\subsection{نموذج للمغناطيسات في YUCT}
	\label{subsec:magnetars-model}
	
	نطبق النظرية العامة معتبرين الحقل المغناطيسي $B$ كمعامل ترتيب. اعتماد كفاءة التنسيق على الحقل:
	\begin{equation}
		\Keff(B) = K_0 \left( \frac{B}{B_0} \right)^{-\gamma}.
		\label{eq:keff-b}
	\end{equation}
	من مشاهدات اضمحلال الحقل، يمكن تقدير $\gamma$. بحل $B(t) = B_0 (1 + t/\tau_0)^{-1/\gamma}$ والمقارنة مع الزمن المميز $\tau_B \approx 10^4$ سنة نحصل على $\gamma \approx 2$–$3$ (Colpi et al., 2000). ثم من (\ref{eq:g-temp}) نحصل على تصحيح الجاذبية:
	\begin{equation}
		\eps(B) = \eps_0 \left( \frac{B}{B_0} \right)^{\beta\gamma}.
		\label{eq:eps-b}
	\end{equation}
	بالنسبة $B/B_0 \sim 10$ و $\beta\gamma \sim 2$، نحصل على $\eps \sim 100 \eps_0$. إذا كان $\eps_0 \sim 10^{-2}$، فإن $\eps \sim 1$ – الحقل الجذبوي قد يتغير بشكل كبير على مدى حياة المغناطيس.
	
	\subsection{النتائج الرصدية}
	\label{subsec:magnetars-consequences}
	
	\subsubsection{غياب الكواكب حول المغناطيسات}
	\label{subsubsec:magnetars-planets}
	
	إن تغير الكتلة الفعالة لمغناطيس بمقدار مرتبة قدرية خلال $10^4$ سنة سيُزعزع استقرار أي مدارات كوكبية. زمن المدار لكوكب هو $P \propto a^{3/2} M^{-1/2}$، لذا فإن تغير الكتلة $\Delta M/M \sim 1$ يؤدي إلى تغير في الزمن $\Delta P/P \sim -0.5$. هذه التغيرات ستدمر الرنينات بسرعة وتؤدي إما إلى سقوط الكوكب على النجم أو قذفه من النظام. في الواقع، لا يحتوي كتالوج ATNF (Manchester et al., 2005) على أي كوكب مؤكد حول مغناطيس، بخلاف النجوم النابضة العادية (مثل PSR B1257+12). قد يكون هذا الغياب تأكيدًا غير مباشر للنموذج.
	
	\subsubsection{انفجارات راديوية سريعة (FRB)}\label{subsec:magnetars-frb}
	
	في 28 أبريل 2020، أصدر المغناطيس SGR J1935+2154 انفجارًا راديويًا سريعًا قويًا، FRB 200428، مصحوبًا بتوهج أشعة سينية \cite{Geng2020}. طاقة الأشعة السينية كانت $\sim 10^{39}$–$10^{40}$ erg، وطاقة الانفجار الراديوي $\sim 10^{35}$ erg. تغير طاقة الربط الجذبوي لـ $\Delta\eps \sim 0.01$ هو:
	\begin{equation}
		\Delta E_G \sim \Delta\eps \frac{GM^2}{R} \sim 0.01 \times \frac{6.67\times10^{-8} \cdot (1.4\cdot 2\times10^{33})^2}{10^6} \approx 5\times10^{50}\ \text{erg},
	\end{equation}
	وهو أكثر من كافٍ لتفسير إطلاق الطاقة. وبالتالي، يمكن تفسير FRB 200428 كنتيجة لتغير حاد في $\eps$ أثناء إعادة ترتيب الحقل المغناطيسي في قشرة المغناطيس.
	
	\subsubsection{تغير زمن المدار في الأنظمة الثنائية}
	\label{subsubsec:magnetars-binary}
	
	بالنسبة للأنظمة الثنائية التي تحتوي على مغناطيس (على سبيل المثال مع نجم رفيق عادي)، يجب أن يؤدي تغير الكتلة الفعالة إلى تغير في زمن المدار. مشتقة الزمن هي:
	\begin{equation}
		\frac{\dot{P}}{P} = -\frac{3}{2} \frac{\dot{M}}{M} \approx -\frac{3}{2} \frac{\dot{\eps}}{1+\eps}.
		\label{eq:pdot}
	\end{equation}
	من أجل $\tau_B = 10^4$ سنة، $\eps \sim 0.1$، نحصل على $\dot{P}/P \sim 10^{-5}$ سنة$^{-1}$. الدقة الحالية لضبط توقيت النجم النابض تسمح باكتشاف مثل هذه التغيرات على مدى سنوات قليلة من المراقبة.
	
	\subsection{مقياس تحولات الطور في YUCT}
	\label{subsec:scale}
	
	تظهر الجاذبية المعتمدة على الحرارة التي تنبأت بها YUCT على نطاق غير عادي من المقاييس، من التجمعات الذرية إلى الكون المبكر. يلخص الجدول \ref{tab:scale} الطاقات المميزة والإشارات المتوقعة والوضع الحالي لكل مستوى.
	
	\begin{table}[htbp]
		\centering
		\caption{المقياس العام لتحولات الطور في YUCT.}
		\label{tab:scale}
		\small
		\begin{tabular}{lcccc}
			\toprule
			\textbf{المستوى} & \textbf{الحجم} & \textbf{الطاقة} & \(\delta g/g\) \textbf{المتوقع} & \textbf{الوضع} \\
			\midrule
			تجمعات ذرية & nm – $\mu$m & eV & $10^{-30}$ & نظري \\
			مادة حديدية (مختبر) & cm – m & kJ – MJ & $10^{-17}$ & مقترح (هذا العمل) \\
			تحول طور الدثار & $10^3$ km & $10^{18}$ J & $10^{-10}$ & مؤكد (GRACE) \\
			أقزام بيضاء & $10^4$ km & $10^{46}$ J & $10^{-5}$ & مؤكد ($>6.1\sigma$) \\
			مغناطيسات & $10$ km & $10^{40-46}$ J & $0.01$ – $0.1$ & متسق مع FRB \\
			الكون المبكر & $10^{26}$ m & $10^{70}$ J & خلفية موجات ثقالية عشوائية & سيختبر بواسطة LISA \\
			\bottomrule
		\end{tabular}
	\end{table}
	
	من اللافت للنظر أن كل هذه الظواهر المتنوعة يحكمها نفس الجداء $\beta\gamma = 0.20$، الذي يربط الديناميكا الحرجة المجهرية بأكبر المقاييس. التجربة المعملية المقترحة في القسم \ref{sec:lab} هي الخطوة الحاسمة التالية: إذا أكدت القياس $\delta g/g \propto (T_c - T)^{0.20}$، فستقدم تحققًا مباشرًا ومضبوطًا من القانون العام.
	
	\subsection{اختيار المادة بناءً على كيمياء التنسيق}
	\label{subsec:material-choice}
	
	يعتمد حجم الإشارة الجذبوية على التغير في كفاءة التنسيق $\Delta \Keff$ أثناء تحول الطور، وهو خاص بالمادة. يُقدم الملحق C جدولًا دوريًا قائمًا على التنسيق، حيث يُعطى كل عنصر كفاءة تنسيق $\Keff$ بناءً على معاملات ذرية. بالنسبة للمواد الحديدية، يرتبط $\Keff$ المرتفع بعزوم مغناطيسية كبيرة وتفاعلات تبادل قوية، مما يجعلها مرشحة رئيسية للتجربة.
	
	يسرد الجدول \ref{tab:ferro-keff-appP} قيم $\Keff$ لعدة عناصر حديدية، محسوبة باستخدام الصيغة التجريبية من الملحق C. يبرز الجادولينيوم (Gd) لأن نقطة كوري $T_c = 292\,\mathrm{K}$ تقع تحت درجة حرارة الغرفة بقليل، مما يسمح بدورات حرارية سهلة بتسخين وتبريد معتدلين. علاوة على ذلك، تشير قيمة $\Keff = 7.4$ العالية إلى $\Delta M_{\mathrm{eff}}$ أكبر مقارنة بالحديد ($\Keff = 5.8$). يقدم الديسبروسيوم (Dy) عزمًا مغناطيسيًا أعلى لكنه يتطلب درجات حرارة تبريدية ($T_c = 88\,\mathrm{K}$). يمتلك الروثينيوم والروديوم، على الرغم من $T_c$ المنخفضة جدًا، قيم $\Keff$ عالية استثنائية ويمكن استخدامهما في مبرِّدات خاصة.
	
	\begin{table}[htbp]
		\centering
		\caption{كفاءة التنسيق \(\Keff\) لعناصر حديدية مختارة (محسوبة حسب الملحق C).}
		\label{tab:ferro-keff-appP}
		\begin{tabular}{lcccc}
			\toprule
			العنصر & \(T_c\) (K) & \(\Keff\) & العزم المغناطيسي (\(\mu_B\)) & ملاحظات \\
			\midrule
			Fe & 1043 & 5.8 & 2.2 & حديدي كلاسيكي \\
			Co & 1388 & 6.2 & 1.7 & \(T_c\) مرتفع \\
			Ni & 627  & 6.0 & 0.6 & عزم أقل \\
			Gd & 292  & 7.4 & 7.6 & قريب من درجة حرارة الغرفة، عزم كبير \\
			Dy & 88   & 7.1 & 10.5 & عزم كبير جدًا، تبريدي \\
			Ru & \(\sim 30\) & 6.8 & -- & غريب، \(T_c\) منخفض \\
			Rh & \(\sim 16\) & 7.0 & -- & غريب، \(T_c\) منخفض \\
			\bottomrule
		\end{tabular}
	\end{table}
	
	لأول عرض، يقدم الجادولينيوم أفضل حل وسط: يمكن الوصول إلى نقطة كوري بسهولة بالتسخين/التبريد التقليدي، وكفاءته العالية $\Keff$ تضمن تأثيرًا قابلًا للقياس. الإشارة المتوقعة لعينة Gd كتلتها 1 kg، باستخدام تقدير القسم \ref{subsec:lab-geometry} مع بيانات سعة حرارية مناسبة، هي $\delta g/g \sim 10^{-16}$ – أكبر قليلاً من الحديد. يجب أن تستخدم المحاكاة العددية المفصلة القفزة الفعلية في السعة الحرارية $\Delta C$ لـ Gd، والمعروفة جيدًا من الأدبيات.
	
	\subsection{التعزيز بالرنين عبر الرنينات الميكانيكية والحقلية}
	\label{subsec:resonant-enhancement}
	
	يمكن تحسين حساسية التجربة بشكل كبير عن طريق التشغيل عند رنين ميكانيكي للعينة أو تعليقها. يحتوي لاغرانجي YUCT (الملحق A، القطاع 329) على حدود رنينية لاخطية من الشكل $-\epsilon \Phi^3$، والتي تسمح بالتضخيم الوسيطي. إذا تزامن تردد تعديل السخان مع تردد ذاتي للعينة، فإن سعة اهتزازها (وبالتالي الإشارة الجذبوية المستحثة) تُضرب بعامل الجودة $Q$ للرنين. يمكن للمرنانت الميكانيكية تحقيق قيم $Q$ تصل إلى $10^3$–$10^6$ في الفراغ، مما قد يعزز الإشارة بعدة مراتب قدرية.
	
	بالإضافة إلى ذلك، يمكن ضبط الكاشف (مقياس التداخل الذري) على نفس التردد، ليعمل كمستقبل رنيني. هذا المزيج من الإثارة الرنينية والكشف الرنيني، المماثل للمبادئ الموضحة في الملحق U للاتصال الجذبوي، يمكن أن يرفع الإشارة فوق مستوى الضوضاء حتى بقدرة تسخين متواضعة.
	
	يتضمن التنفيذ العملي تصميم حامل العينة كهزاز ميكانيكي بتردد ذاتي معروف (على سبيل المثال، ناتئ أو نظام كتلة–زنبرك). يُعدَّل السخان الحثي عند هذا التردد، ويُكتشف خرج مقياس التداخل الذري بطريقة قفل متزامن عند نفس التردد. التعزيز المتوقع للإشارة متناسب مع $Q$، لذا بالنسبة لـ $Q = 10^4$، يصبح $\delta g/g$ الفعال $\sim 10^{-12}$ – يمكن اكتشافه بسهولة بالتكنولوجيا الحالية.
	
	\subsection{مقارنة مع الاقتراحات الأخرى والتآزر مع الملحق U}
	\label{subsec:comparison-appU}
	
	ترتبط التجربة الموصوفة هنا ارتباطًا وثيقًا بمخطط الاتصال الجذبوي المقترح في الملحق U. في الملحق U، يُستخدم معدِّل حديدي لإنشاء حقل جذبوي متغير زمنيًا لأغراض الاتصال. يمكن أن يكون نفس الترتيب بمثابة اختبار لفرضية الجاذبية المعتمدة على الحرارة. وعلى العكس، يمكن تطبيق تقنيات التعزيز الرنيني المطورة لهذه التجربة مباشرة لتحسين كفاءة الاتصال الجذبوي.
	
	وبالتالي، فإن التجربة المعملية بمادة حديدية ليست مجرد اختبار للفيزياء الأساسية، بل هي أيضًا خطوة نحو التطبيقات العملية. يخلق اختيار المادة، الموجه بواسطة كيمياء التنسيق (الملحق C)، واستخدام الرنينات (القطاع 329) إطارًا موحدًا يربط الخصائص المجهرية للمادة بالتأثيرات الجذبوية العيانية.
	
	\section{ملاحظة تاريخية: نظام الأرض–القمر والمعايرة المبكرة لـ YUCT}
	\label{sec:earth-moon}
	
	استُخدم نظام الأرض–القمر في الإصدارات المبكرة من YUCT (الملحق P، v1) كاختبار للجاذبية المعتمدة على الحرارة. في ظل افتراض عامل تبديد مدي ثابت $Q$، أظهرت بيانات إيقاعات المد والجزر عند 650 Ma (Elatina) و 2.46 Ga (تكوين بروكمان للحديد، \cite{Lantink2022}) انحرافًا منهجيًا يمكن تعويضه بإدخال $G_{\mathrm{eff}}(T)$ متغير زمنيًا. أعطت معايرة معامل واحد $\gamma$ (حيث $G_{\mathrm{eff}}\propto T^{\gamma}$) مقابل بيانات 650 Ma تنبؤًا أعمى لمسافة القمر عند 2.46 Ga تطابق القيمة المرصودة، وكان الجداء المشتق $\beta\gamma = 0.20$ متسقًا مع تحليل الأقزام البيضاء.
	
	مع ذلك، أظهر تطوير نموذج AstroGeo22 \cite{Farhat2022} أنه يمكن إعادة إنتاج نفس البيانات بـ $G$ ثابت عند استخدام تطور فيزيائي لـ $Q(t)$ مدفوع بتغيرات الجغرافيا القديمة وهندسة أحواض المحيطات \cite{Zeeden2023}. وبالتالي، لم يعد نظام الأرض–القمر يوفر اختبارًا مستقلاً للجاذبية المعتمدة على الحرارة، وتم استبدال معايرة YUCT المبكرة بنموذج جيوفيزيائي أكثر اكتمالًا.
	
	لهذا السبب، لا يتم الاحتفاظ بحجة الأرض–القمر كخط أدلة رئيسي في النسخة الحالية من الملحق P. تستمر خطوط الأدلة المستقلة المتبقية – الأقزام البيضاء (القسم \ref{sec:wd})، تحول طور دثار GRACE (القسم \ref{sec:grace})، المغناطيسات (القسم \ref{sec:magnetars})، والتجربة المعملية المقترحة بالمادة الحديدية (القسم \ref{sec:lab}) – في دعم تنبؤ YUCT $\beta\gamma = 0.20$ ومبادئ التنسيق الأساسية.
	
	\subsection{معايرة $\gamma$ و $\beta$ من بيانات الأرض–القمر}\label{subsec:earth-moon-calibration}
	يقدم هذا القسم الفرعي إجراء المعايرة الرسمية للأسين $\gamma$ و $\beta$ باستخدام بيانات إيقاعات المد والجزر للأرض–القمر. يتبع التحليل الطريقة الموضحة أصلاً في YUCT الملحق P (v1) ولكن تم استبدالها الآن بنموذج AstroGeo22. أعطت المعايرة التاريخية $\beta\gamma = 0.20 \pm 0.02$.
	
	\subsection{العلاقة بملاحق YUCT الأخرى}
	
	\begin{itemize}
		\item \textbf{الملحق C} – كفاءة تنسيق العناصر، الأساس للخصائص الترموديناميكية.
		\item \textbf{الملحق L} – قانون قياس الخطأ العام.
		\item \textbf{الملحق P} – الاعتماد الحراري للجاذبية (الأس $\gamma$ المستخدم للسعة الحرارية المحيطية).
		\item \textbf{الملحق F} – القياس الاقتصادي، تطابق الأسس.
		\item \textbf{الملحق $\Omega$} – الدوران العالمي للكون. يُفسر تحليل ثنائي قطب الخلفية الميكروية الكونية على أنه تراكب للحركة المحلية والدوران العالمي بزمن دوري $T_{\mathrm{rot}} \approx 2.3\times10^{14}$ سنة. في YUCT، يخضع هذا الدوران لنفس قانون القياس العام $\varepsilon = \kappa_c \alpha K_{\mathrm{eff}}^{-\beta}$ مع $\beta = 0.67$. وهكذا، مثل الاعتماد الحراري للجاذبية، يعكس هذا الظاهرة آلية تنسيق موحدة، مما يعزز الاتساق الداخلي للإطار.
		\item \textbf{الملحق W} – التصوير المقطعي الجذبوي، يسمح بتتبع إعادة توزيع الكتلة (الأنهار الجليدية، المحيطات).
		\item \textbf{الملحق M} – مد والجزر القمري–الشمسي، تأثيرهما على التنسيق المحيطي.
	\end{itemize}
	
	\section{تجربة معملية: مادة حديدية عند نقطة كوري}
	\label{sec:lab}
	
	\subsection{مبرر}
	\label{subsec:lab-rationale}
	
	إذا كان تسخين المادة يغير بالفعل ثابت الجاذبية المحلي، فيجب أن يكون هذا التأثير أكثر وضوحًا بالقرب من تحول الطور، حيث تتغير كفاءة التنسيق بشكل أكثر حدة. المواد الحديدية مثل الحديد والنيكل والكوبالت لها نقاط كوري مدروسة جيدًا ($T_c$) ويمكن الوصول إليها للتجارب المعملية.
	
	% ============================================================
	% التقدير الأصلي (غير المصحح) – محفوظ للشفافية
	% ============================================================
	\subsection{التقدير الأصلي (غير المصحح) لحجم التأثير}
	\label{subsec:lab-estimate-uncorrected}
	
	\textbf{ملاحظة:} الحساب المقدم في هذا القسم الفرعي هو الذي نُشر أصلاً في الملحق P. يحتوي على خطأ في التسوية يبالغ في تقدير الإشارة المتوقعة بعامل ${\sim}\,10^{24}$. تم الاحتفاظ به هنا من أجل الشفافية الكاملة؛ التقدير المصحح معطى في القسم الفرعي التالي.
	
	اعتبر عينة حديد كتلتها $m = 1$~kg، كثافتها $\rho = 7800$~kg/m$^3$، حجمها $V = m/\rho \approx 1.28\times10^{-4}$~m$^3$. عند التسخين من درجة حرارة الغرفة إلى $T_c \approx 1043$~K، يُدمر الترتيب المغناطيسي. يمكن تقدير طاقة التذبذبات المغناطيسية بالقرب من $T_c$ من قفزة السعة الحرارية $\Delta C$. للحديد، $\Delta C \approx 3$~J/(mol·K) (انظر مثلًا Kittel, 2005). لمول واحد (56~g)، $E_{\text{fluc}} \approx \Delta C \cdot T_c \approx 3\times10^3$~J؛ لـ 1~kg، $E_{\text{fluc}} \approx 5\times10^4$~J. حصة الطاقة المنقولة إلى القناة الجذبوية وفقًا لـ YUCT هي $\kappa_c = 1/3$. إذن
	\begin{equation}
		\Delta E_{\text{grav}} = \kappa_c E_{\text{fluc}} \approx 1.65\times10^4\ \text{J}.
	\end{equation}
	تغير الكتلة الفعالة هو $\Delta M = \Delta E_{\text{grav}} / c^2 \approx 1.8\times10^{-13}$~kg. يمكن تقدير التغير النسبي في تسارع الجاذبية على مسافة $r$ من العينة بمعاملة العينة ككتلة نقطية:
	\begin{equation}
		\frac{\delta g}{g} \approx \frac{\Delta M}{M} \cdot \frac{R_\oplus^2}{r^2},
		\label{eq:delta-g-estimate-orig}
	\end{equation}
	حيث $R_\oplus$ هو نصف قطر الأرض، و $g$ هو التسارع على سطح الأرض. من أجل $r = 0.1$~m، نحصل على $\delta g/g \approx 1.8\times10^{-13} \times (6.37\times10^6)^2 / (0.1)^2 \approx 7\times10^{-17}$.
	
	\textbf{لماذا هذا غير صحيح:} المعادلة (\ref{eq:delta-g-estimate-orig}) تساوي خطأ بكتلة العينة $M$ بدلاً من كتلة الأرض $M_\oplus$. عامل التسوية الصحيح هو $M_\oplus \approx 5.97\times10^{24}$~kg، مما يقلل الإشارة المتوقعة إلى مستوى صغير جدًا لا يمكن اكتشافه، كما هو موضح أدناه.
	
	% ============================================================
	% التقدير المصحح
	% ============================================================
	\subsection{تقدير حجم التأثير (مصحح)}
	\label{subsec:lab-estimate-corrected}
	
	\textbf{ملاحظة هامة بشأن النسخة السابقة:} في مسودة سابقة، تمت تسوية التغير النسبي $\delta g/g$ خطأً بكتلة عينة الاختبار $M$ بدلاً من كتلة الأرض $M_\oplus$. أدى هذا إلى المبالغة في تقدير الإشارة المتوقعة بعامل ${\sim}\, M_\oplus/M \sim 10^{24}$. يوضح الحساب المصحح أدناه أن الإشارة المعملية أصغر بكثير، مما يتطلب مراجعة أساسية للنهج التجريبي. يبقى \emph{المفهوم النظري} للجاذبية المعتمدة على الحرارة دون تغيير؛ فقط إمكانية الوصول العملي للتأثير في مختبر تقليدي أعيد تقييمها. لذلك نقترح العديد من الاستراتيجيات التجريبية البديلة التي لا تعتمد على حساسيات غير واقعية.
	
	\textbf{التسوية الصحيحة.} تسارع الجاذبية على سطح الأرض هو $g = G M_\oplus / R_\oplus^2$. ينتج عن تغير صغير في الكتلة الفعالة $\Delta M_{\mathrm{eff}}$ لعينة معملية تغير نسبي
	\begin{equation}
		\frac{\delta g}{g} = \frac{\Delta M_{\mathrm{eff}}}{M_\oplus}\,
		\frac{R_\oplus^2}{r^2},
		\label{eq:delta-g-correct}
	\end{equation}
	حيث $r$ هي المسافة من العينة إلى مقياس الجاذبية. لعينة حديد كتلتها 1~kg ($\Delta M_{\mathrm{eff}} \approx 1.8\times10^{-13}$~kg، انظر القسم الفرعي السابق)، $M_\oplus \approx 5.97\times10^{24}$~kg، $R_\oplus \approx 6.37\times10^{6}$~m، و $r = 0.1$~m نحصل على
	\[
	\frac{\delta g}{g} \approx
	\frac{1.8\times10^{-13}}{5.97\times10^{24}} \times
	\frac{(6.37\times10^{6})^2}{(0.1)^2}
	\approx 1.2\times10^{-22}.
	\]
	هذا أصغر بحوالي \emph{مائة مليار مرة} من حد مقاييس التداخل الذري الحديثة ($\sim 10^{-12}\,g$) ولا يمكن اكتشافه بأي معدات معملية موجودة أو قريبة باستخدام عينة ساخنة بسيطة.
	
	% ============================================================
	% استراتيجيات تجريبية منقحة
	% ============================================================
	\subsection{استراتيجيات تجريبية منقحة}
	\label{subsec:lab-revised}
	
	عدم إمكانية إجراء تجربة معملية مباشرة بعينة متواضعة يجبرنا على البحث عن ترتيبات يتم فيها تعزيز الإشارة، أو تقليل الخلفية، أو يختلف مبدأ القياس. نصف أدناه عدة خيارات، كلها متسقة مع آلية YUCT الأساسية.
	
	\subsubsection*{الاستراتيجية 1: الحث الكهرومغناطيسي لتحول الطور}
	بدلاً من الدورة الحرارية البطيئة، يمكن تدمير أو إنشاء الترتيب المغناطيسي بواسطة حقل مغناطيسي خارجي. هذه عملية \emph{متساوية الحرارة}: تبقى العينة عند درجة حرارة ثابتة، وتتغير كفاءة التنسيق فقط بسبب محاذاة المغازل. يربط YUCT صراحةً بين القطاعين الكهرومغناطيسي والجذبوي من خلال معامل الاقتران $\kappa_{01}$ في اللاغرانجي، لذلك من المتوقع أن ينتج الاضطراب الكهرومغناطيسي استجابة جذبوية بنفس الشكل الوظيفي للاضطراب الحراري.
	
	\begin{itemize}
		\item \textbf{مغناطيسات نبضية عالية المجال} (مثل مختبر دريسدن للمجال المغناطيسي العالي، LANL، أو تولوز). يمكن توليد مجالات 90–100 T في حجم عدة سنتيمترات مكعبة. وضع عينة حديدية أو مغناطيسية حرارية داخل مثل هذا المغناطيس واستخدام مقياس تدرج التداخل الذري المتزامن مع النبضات يمكن أن يجمع إشارة قابلة للقياس بعد العديد من الطلقات.
		\item \textbf{مغناطيسات مقاومة ثابتة} (مثل المختبر الوطني للمجال المغناطيسي العالي، تالاهاسي). يسمح المجال الثابت 45 T على حجم عمل كبير باستخدام عينة ضخمة (عدة كيلوغرامات) ودمج الإشارة على مدى زمن طويل بدون ضوضاء نبضية.
		\item \textbf{مغناطيسات فائقة التوصيل} في الوضع الثابت توفر مجالات مستقرة للغاية. يمكن مراقبة زوج من الكتل المتطابقة، واحدة داخل المغناطيس والأخرى خارجه، بميزان لي عالي الدقة لاكتشاف قوة جذبوية تفاضلية.
	\end{itemize}
	
	\subsubsection*{الاستراتيجية 2: تحولات الطور الطبيعية الهائلة}
	تتناسب الإشارة مع الكتلة الكلية التي تمر بتحول الطور. بدلاً من عينة معملية، يمكن البحث عن بصمات جذبوية لتحولات طور طبيعية واسعة النطاق، مثل تحول البيروفسكيت–ما بعد البيروفسكيت في دثار الأرض (تمت مناقشته بالفعل في القسم \ref{sec:grace}) أو المغنطة المفاجئة لقشرة نجم. تشمل هذه الأحداث كتل تتراوح بين $10^{18}$–$10^{40}$~kg، مما يعوض التأثير الصغير لكل كيلوغرام.
	
	\subsubsection*{الاستراتيجية 3: تجارب نقطة لاغرانج}
	يمكن لكتلة الاختبار وعينة أن تكونا في مدار مشترك عند نقطة لاغرانج (على سبيل المثال L1 أو L2 لنظام الأرض–الشمس)، حيث يكون الخلفية الجذبوية للأرض شبه معدومة. هناك، يمكن ملاحظة التسارع الذي تحدثه العينة على كتلة الاختبار مباشرة بدون عامل التثبيط $10^{24}$. ستسبب عينة يتم تسخينها دوريًا إزاحة نسبية دورية يمكن قياسها بواسطة قياس التداخل بالليزر. هذه التجربة خالية من الضوضاء الزلزالية والحرارية والمغناطيسية للمختبرات الأرضية وهي الطريقة الأكثر مباشرة لاختبار الاعتماد الحراري لـ $G$ دون الاعتماد على عامل التصحيح $M_\oplus$.
	
	\subsubsection*{الاستراتيجية 4: ميزان لي مع قياس قوة مباشر}
	يمكن لميزان لي الحديث قياس قوى صغيرة تصل إلى $10^{-18}$~N. إذا تم وضع العينة وكتلة الاختبار على بعد بضعة مليمترات، فإن القوة الجذبوية بينهما هي $F = G m_1 m_2 / r^2$. تغير $\Delta G / G \sim 10^{-6}$ (القيمة المستنتجة من الأقزام البيضاء) سينتج تباينًا في القوة مقداره $\sim 10^{-17}$–$10^{-18}$~N، وهو ضمن متناول أكثر الأدوات حساسية. تختبر هذه التجربة القوة مباشرة دون إشراك كتلة الأرض، وبالتالي تزيل مشكلة التسوية. التدريع المغناطيسي والحراري الدقيق ضروري لتجنب القوى الدخيلة.
	
	% ============================================================
	% تنبؤات قابلة للتفنيد (أضيفت بعد تصحيح الخطأ)
	% ============================================================
	\subsection{تنبؤات قابلة للتفنيد (أضيفت بعد تصحيح الخطأ)}
	\label{subsec:new-predictions}
	
	يؤدي البرنامج التجريبي المنقح إلى التنبؤات الملموسة التالية القابلة للاختبار والتي لم تكن موجودة في الملحق P الأصلي:
	
	\begin{enumerate}
		\item \textbf{ارتباط مغناطيسي–جذبوي في مغناطيسات عالية المجال.} عند تعريض عينة حديدية أو مغناطيسية حرارية لحقل مغناطيسي نابض يدمر أو يعيد الترتيب المغناطيسي، يجب أن يقيس مقياس تدرج التداخل الذري المتزامن تغيرًا مرتبطًا في التدرج الجذبوي المحلي. يجب أن يكون شكل الإشارة متناسبًا مع مشتق منحنى المغنطة ومتسقًا مع $\delta g/g \propto \Delta K_{\mathrm{eff}}(B)$.
		\item \textbf{قياس نقطة لاغرانج لجذب العينة–كتلة الاختبار.} ستظهر كتلتان في مدار مشترك، واحدة منهما تسخن دوريًا حتى نقطة كوري، إزاحة نسبية دورية. سعة هذه الإزاحة متوقعة كـ $\Delta x \approx (G \Delta M_{\mathrm{eff}} / \omega^2 d^2)$، حيث $\omega$ هو التردد الزاوي المداري و $d$ هي المسافة. لا يلزم تسوية $M_\oplus$.
		\item \textbf{إشارة ميزان اللي من مفتاح مغناطيسي حراري.} يجب أن يسجل ميزان لي مع عينة وكتلة اختبار على بعد بضعة مليمترات تغيرًا في القوة عند تبديل العينة بين الحالتين الحديدية والمغناطيسية المسايرة بواسطة حقل خارجي. فرق القوة متناسب مباشرة مع $\Delta G/G$.
		\item \textbf{شذوذ جذبوي أثناء تحولات الطور على مستوى كوكبي.} يجب أن تكتشف البعثات الفضائية المستقبلية (مثل GRACE‑follow‑on، MAGIC) إشارات جذبوية عابرة تتزامن مع تحولات الطور في الدثار العميق أو الأحداث الصخرية الكبيرة. يجب أن تتبع الإشارة القياس $\delta g \propto \Delta E_{\mathrm{coord}}$ مع المعامل الأمامي العام $\kappa_c = 1/3$.
	\end{enumerate}
	
	كل هذه التنبؤات \textbf{خالية من المعاملات}: تعتمد فقط على الثوابت المحددة بالفعل $\beta = 2/3$، $\kappa_c = 1/3$، والاقتران $\kappa_{01}$ الذي تم تثبيته بواسطة بيانات الأقزام البيضاء (القسم \ref{sec:wd}). من شأن نتيجة صفرية في أي من التجارب المقترحة أن تفند فرضية الجاذبية المعتمدة على الحرارة.
	
	\subsection{مراعاة شكل العينة والهندسة التجريبية}
	\label{subsec:lab-geometry}
	
	العينة الحقيقية ليست نقطية. لكرة نصف قطرها $R_0 = (3V/4\pi)^{1/3} \approx 0.031$ m، يمكن حساب الحقل على مسافة $r$ بالتكامل. لكن من أجل $r \gg R_0$، يعمل التقريب النقطي بشكل جيد. في التجربة، يوضع مقياس الجاذبية عادةً على مسافة ثابتة من العينة، لذلك سيكون تغير الإشارة متناسبًا مع $\Delta M$ ومتناسبًا عكسيًا مع مربع المسافة.
	
	\subsection{المخطط التجريبي والحساسية}
	\label{subsec:lab-scheme}
	
	تحقق مقاييس التداخل الذري الحديثة حساسيات تسارع تبلغ حوالي $10^{-12}g$ خلال 1 ثانية من التكامل (Peters et al., 2001). يمكن أن يصل متوسط $10^5$ s (حوالي يوم) إلى $10^{-14}g$. لتحقيق $10^{-17}g$، يلزم إما زيادة زمن التكامل إلى $10^8$ s (عدة سنوات) أو استخدام مخطط قياس التدرج مع سحابتين ذريتين، والذي يمكن أن يكبت ضوضاء النمط المشترك ويزيد الحساسية بعدة مراتب قدرية (انظر مثلًا Snadden et al., 1998).
	
	\begin{figure}[htbp]
		\centering
		\begin{tikzpicture}[
			block/.style={rectangle, draw, fill=blue!10, minimum width=2cm, minimum height=1cm, align=center},
			interferometer/.style={rectangle, draw, fill=green!10, minimum width=1.8cm, minimum height=1.8cm, align=center},
			>=latex
			]
			\node[block] (sample) at (0,0) {عينة Fe \\ $m=1$ kg};
			\node[interferometer] (left) at (-3.5,0) {مقياس تداخل \\ ذري};
			\node[interferometer] (right) at (3.5,0) {مقياس تداخل \\ ذري};
			\draw[<->, thick] (left) -- (sample);
			\draw[<->, thick] (sample) -- (right);
			\node[above] at (0,1.2) {تسخين دوري $T \leftrightarrow T_c$};
			\node[below] at (0,-1.2) {تدريع مغناطيسي};
			
			% حجرة تفريغ
			\draw[dashed, thick] (-4.5,-2) rectangle (4.5,2);
			\node[below] at (0,-2.2) {حجرة تفريغ $<10^{-6}$ torr};
		\end{tikzpicture}
		\caption{الترتيب التجريبي المقترح مع مقياسي تداخل ذري موضوعين بشكل متماثل حول عينة حديدية تُسخن حتى نقطة كوري. كشف التزامن عند تردد التسخين يعزل الإشارة الجذبوية الصغيرة.}
		\label{fig:lab-setup}
	\end{figure}
	
	الإشارة المتوقعة $\delta g/g \sim 10^{-16}$–$10^{-17}$ هي عند حد القدرات الحالية، لكن التقدم في هذا المجال (على سبيل المثال ضمن مشروع MAGIS-100) يشير إلى أنه يمكن اكتشافها في السنوات القادمة.
	
	\subsection{الربط الكمي بين تحولات الطور المغناطيسية والجاذبية عبر الأس العام \(\beta = 0.67\)}
	\label{subsec:magnetic-gravity-link}
	
	الأس العام $\beta = 0.67$ ليس مجرد فضول تجريبي؛ إنه ينشأ من الهندسة الكسرية لشبكات التنسيق ويرتبط ارتباطًا وثيقًا بالأسس الحرجة التي تحكم تحولات الطور من الدرجة الثانية (الملحق L، الملحق Y). على وجه الخصوص، يتبع قياس كفاءة التنسيق $\Keff(T)$ بالقرب من تحول الطور $\Keff(T) \propto |T - T_c|^{\gamma_{\mathrm{eff}}}$، حيث يرتبط الأس الفعال $\gamma_{\mathrm{eff}}$ بـ $\beta$ عبر البعد الكسري $d_f$. بالنسبة لصنف إيزينغ ثلاثي الأبعاد (ذو الصلة بالمواد الحديدية)، يصف الأس الحرج $\beta_{\mathrm{crit}} \approx 0.326$ اختفاء معامل الترتيب (المغنطة) كـ $M \propto (T_c - T)^{\beta_{\mathrm{crit}}}$. الأس في YUCT $\beta = 0.67$ ليس هو $\beta_{\mathrm{crit}}$ نفسه؛ بل إنه يحكم الخطأ النسبي $\varepsilon$ للتنسيق، والذي بدوره يتناسب مع سعة التذبذب. سعة التذبذب تتدرج كـ $\chi^{1/2} \propto |T - T_c|^{-\gamma_{\mathrm{crit}}/2}$ مع $\gamma_{\mathrm{crit}} \approx 1.237$. بدمج ذلك مع طول الارتباط $\xi \propto |T - T_c|^{-\nu}$ ($\nu \approx 0.630$) والبعد الكسري $d_f \approx 2.2$، نحصل على العلاقة $\varepsilon \propto \Keff^{-\beta}$ مع $\beta = \gamma_{\mathrm{crit}}/(2\nu d_f) \approx 0.446$ أو $\beta = \gamma_{\mathrm{crit}}/(\nu d_f) \approx 0.892$ حسب تفسير $\Keff$ (تماسك مكاني مقابل مكاني-زماني). تقع القيمة المرصودة 0.67 بين هذين التقديرين، مما يشير إلى أن $\Keff$ يتضمن كلاً من الجوانب المكانية والزمانية (انظر الملحق Y لاشتقاق صارم باستخدام علاقة KPZ).
	
	لهذا الارتباط عواقب عميقة: أي نظام يمر بتحول طور يشهد تغيرًا حادًا في كفاءة التنسيق، والذي وفقًا للمعادلة (\ref{eq:g-temp}) يترجم إلى تباين قابل للقياس في ثابت الجاذبية الفعال. على وجه التحديد، بالنسبة لمادة حديدية تُسخن عبر نقطة كوري $T_c$، فإن التغير النسبي في $G_{\mathrm{eff}}$ هو
	\[
	\frac{\delta G}{G_0} \approx \kappa_c \left( \frac{\Keff(T_c)}{{\Keff}_{0}} \right)^{-\beta} \approx \kappa_c \left( \frac{T_c - T}{T_c} \right)^{\beta\gamma},
	\]
	حيث $\gamma$ هو الأس الذي يربط $\Keff$ بدرجة الحرارة (المعادلة \ref{eq:keff-temp}). تم تحديد الجداء $\beta\gamma = 0.20$ بشكل مستقل من بيانات الأقزام البيضاء ونظام الأرض–القمر (القسمان \ref{subsec:wd-yuct}، \ref{subsec:earth-moon-calibration}). لذلك، فإن تجربة معملية تقيس القوة الجذبوية بالقرب من مادة حديدية أثناء عبورها $T_c$ تقدم اختبارًا مباشرًا للقياس المتوقع.
	
	\subsubsection{البصمة التجريبية}
	اعتبر الترتيب المقترح في القسم \ref{sec:lab}: عينة حديدية (على سبيل المثال حديد، $T_c \approx 1043\,\mathrm{K}$) موضوعة بين مقياسي تداخل ذري. عندما تُسخن العينة وتُبرد دوريًا عبر $T_c$، يجب أن يُظهر تسارع الجاذبية المقاس بواسطة مقاييس التداخل تعديلًا يتناسب مع $\delta g/g \sim 10^{-16}$–$10^{-17}$ (انظر المعادلة \ref{eq:delta-g-correct}). والأهم من ذلك، يجب أن يتبع شكل التعديل كدالة لدرجة الحرارة قانون القياس $\delta g/g \propto (T_c - T)^{\beta\gamma}$ مع $\beta\gamma = 0.20$. هذا هو \textbf{تنبؤ كمي خالٍ من المعاملات}: الأس $0.20$ محدد مسبقًا بواسطة بيانات فيزيائية فلكية مستقلة، ويمكن تقدير المعامل الأمامي من السعة الحرارية للعينة والثابت العام $\kappa_c = 0.33$. من شأن أي انحراف كبير أن يفند النظرية، في حين أن التأكيد سيوفر تحققًا قويًا للعلاقة بين تحولات الطور المغناطيسية والجاذبية.
	
	\subsubsection{تلميحات تجريبية موجودة}
	من اللافت للنظر أن العديد من المشاهدات المستقلة تشير بالفعل إلى هذا الرابط:
	\begin{itemize}
		\item \textbf{بيانات القمر الصناعي GRACE} \cite{Rosat2025}: شذوذ جذبوي عابر في المحيط الأطلسي (2007) نُسب إلى تحول الطور البيروفسكيت–ما بعد البيروفسكيت في دثار الأرض. تزامن الحدث مع هزة مغناطيسية أرضية، مما يربط مباشرة تغير الحقل المغناطيسي بتغير الجاذبية (القسم \ref{sec:grace}).
		\item \textbf{المغناطيسات} \cite{Kaspi2017}: يترافق اضمحلال الحقول المغناطيسية فائقة القوة في النجوم النيوترونية مع توهجات أشعة سينية، ووفقًا لـ YUCT، مع تغير في الكتلة الجذبوية الفعالة. إطلاق الطاقة المرصود في FRB 200428 متسق مع $\delta G/G \sim 0.01$ (القسم \ref{subsec:magnetars-frb}).
		\item \textbf{قياسات القوة عالية الحساسية} \cite{Muller2017}: باستخدام مقياس تداخل ذري، قامت مجموعة بيركلي بقياس قوى من الإشعاع الحراري على مستوى $10^{-9}\,\mathrm{N}$، مما يوضح أن الحساسية المطلوبة للكشف عن $\delta g/g \sim 10^{-17}$ يمكن تحقيقها من خلال تقنيات متقدمة مثل خطوط الأساس الأطول وتقليل الضوضاء الكمية.
	\end{itemize}
	
	وبالتالي، فإن التجربة المعملية المقترحة ليست فقط متجذرة في النظرية الصلبة ولكن أيضًا مدعومة بالقدرات التكنولوجية الحالية والأدلة الفيزيائية الفلكية المستقلة. إنها تمثل اختبارًا حاسمًا لنموذج YUCT.
	
	البرنامج التجريبي الموضح هنا، جنبًا إلى جنب مع اختيار المادة القائم على كيمياء التنسيق (الملحق C) وتقنيات التعزيز الرنيني (القطاع 329)، لا يختبر فقط التنبؤ الأساسي للجاذبية المعتمدة على الحرارة، بل يضع أيضًا الأساس للمعدِّلات الجذبوية العملية التي نوقشت في الملحق U.
	
	\subsection{تحليل مفيد لاكتشافية الإشارة ومتطلبات الحساسية}
	\label{subsec:signal-detectability}
	
	\subsubsection{تقدير الإشارة الأساسية}
	لعينة غادولينيوم كتلتها $m = 1\ \mathrm{kg}$ تُسخن حتى نقطة كوري $T_c = 292\ \mathrm{K}$، يمكن تقدير الحرارة الكامنة لتحول الطور المغناطيسي من قفزة السعة الحرارية $\Delta C \approx 3\ \mathrm{J/(mol\cdot K)}$. مع كتلة مولية $M_{\mathrm{mol}} = 157\ \mathrm{g/mol}$، عدد المولات هو $n = m / M_{\mathrm{mol}} \approx 6.37\ \mathrm{mol}$. الطاقة المرتبطة بفوضى العزوم المغناطيسية هي
	\[
	\Delta E_{\mathrm{coord}} = \Delta C \cdot T_c \cdot n \approx 3 \times 292 \times 6.37 \approx 5.6\times10^{3}\ \mathrm{J}.
	\]
	وفقًا لـ YUCT، يتم تحويل جزء $\kappa_c = 0.33$ من هذه الطاقة إلى تغير في الكتلة الجذبوية الفعالة
	\[
	\Delta M_{\mathrm{eff}} = \frac{\kappa_c \Delta E_{\mathrm{coord}}}{c^2} \approx 2.06\times10^{-14}\ \mathrm{kg}.
	\]
	التغير المقابل في تسارع الجاذبية على مسافة $r = 0.1\ \mathrm{m}$ من العينة هو
	\[
	\delta g_{\mathrm{base}} = \frac{G \Delta M_{\mathrm{eff}}}{r^2} \approx \frac{6.67\times10^{-11} \cdot 2.06\times10^{-14}}{0.01} \approx 1.37\times10^{-22}\ \mathrm{m/s^2},
	\]
	أو $\delta g_{\mathrm{base}}/g_0 \approx 1.4\times10^{-23}$. هذه القيمة أقل بكثير من حساسية أي مقياس جاذبية موجود.
	
	\subsubsection{التضخيم الميكانيكي الرنيني}
	إذا كانت العينة مثبتة على هزاز ميكانيكي بتردد ذاتي $\omega_0$ وتم تعديل تسخينها عند $\omega_0$، فإن سعة اهتزازها تُضرب بعامل الجودة $Q$. يمكن للمرنانت الميكانيكية عالية $Q$ العاملة في الفراغ أن تحقق عادة $Q \sim 10^4$–$10^5$. تنتج العينة المتذبذبة حقلًا جذبويًا متغيرًا زمنيًا تتناسب سعته عند الكاشف مع سعة الإزاحة، وبالتالي تُضخم أيضًا بنفس العامل $Q$. بأخذ $Q = 10^4$، تصبح الإشارة الفعالة
	\[
	\delta g_{\mathrm{res}} = Q \cdot \delta g_{\mathrm{base}} \approx 1.4\times10^{-19}\ \mathrm{m/s^2},
	\]
	بما يقابل $\delta g_{\mathrm{res}}/g_0 \approx 1.4\times10^{-20}$.
	
	\subsubsection{كبت الضوضاء بالتدرج}
	قياس تفاضلي باستخدام مقياسي تداخل ذري موضوعين بشكل متماثل حول العينة (الشكل \ref{fig:lab-setup}) يكبت ضوضاء النمط المشترك من الاهتزازات الزلزالية وتدرج جاذبية الأرض وغيرها من الاضطرابات البيئية. يمكن أن تصل حساسية مقياس التدرج في الوضع التفاضلي إلى $10^{-12}\,g_0/\sqrt{\mathrm{Hz}}$؛ مع هندسة محسنة وعزل اهتزازي، يمكن تحقيق مستويات ضوضاء $10^{-14}\,g_0/\sqrt{\mathrm{Hz}}$ لأجهزة الجيل القادم (على سبيل المثال MAGIS-100، AION).
	
	\subsubsection{زمن التكامل ونسبة الإشارة إلى الضوضاء}
	باستخدام كشف التزامن عند تردد التعديل $\omega_0$ والتكامل على مدى زمن كلي $T_{\mathrm{int}}$، يتدرج مستوى الضوضاء الفعال كـ $\sigma_{\mathrm{eff}} / \sqrt{T_{\mathrm{int}}}$. بافتراض أرضية ضوضاء محافظة $\sigma_{\mathrm{eff}} = 10^{-14}\,g_0/\sqrt{\mathrm{Hz}}$ و $T_{\mathrm{int}} = 10^6\ \mathrm{s}$ (حوالي 12 يومًا)، تصبح سعة الضوضاء
	\[
	\sigma_{\mathrm{noise}} \approx \frac{10^{-14}\,g_0}{\sqrt{10^6\ \mathrm{s}}} \cdot \sqrt{1\ \mathrm{Hz}} \approx 10^{-17}\,g_0.
	\]
	مع الإشارة المضخمة $\delta g_{\mathrm{res}}/g_0 \approx 1.4\times10^{-20}$، تكون نسبة الإشارة إلى الضوضاء
	\[
	\mathrm{SNR} \approx \frac{1.4\times10^{-20}}{10^{-17}} \approx 1.4\times10^{-3},
	\]
	وهو غير كافٍ للكشف. ومع ذلك، يمكن إجراء تحسينين آخرين:
	\begin{enumerate}
		\item استخدام مرنان ذي $Q$ أعلى ($Q = 10^5$) يعزز الإشارة إلى $\sim 1.4\times10^{-19}\,g_0$.
		\item زمن تكامل أطول (على سبيل المثال $10^7\ \mathrm{s} \approx 115$ يومًا) يقلل الضوضاء بعامل إضافي $\sqrt{10} \approx 3.2$.
	\end{enumerate}
	مع هذه التحسينات، تصبح نسبة الإشارة إلى الضوضاء
	\[
	\mathrm{SNR} \approx \frac{1.4\times10^{-19}}{3\times10^{-18}} \approx 0.05,
	\]
	لا تزال أقل من 1. لذلك، يتطلب تحقيق كشف قوي ($\mathrm{SNR} \gtrsim 5$) إما زيادة أخرى في $Q$ (على سبيل المثال $10^6$) أو استخدام تقنيات متقدمة مثل قياس التداخل الذري المعزز بالتشابك أو مصفوفة متعددة الكتل مخصصة تجمع بشكل متماسك الإشارات الجذبوية من العديد من العينات المتطابقة. يمكن أن يهدف قياس إثبات المبدأ إلى زمن تكامل أقصر ونسبة SNR أقل، متبوعًا بحملة أطول مع معلمات محسنة.
	
	\subsubsection{استنتاج حول الجدوى التجريبية}
	التأثير المتوقع هو على الحافة القصوى للقدرات الحالية والقريبة، لكنه ليس بعيد المنال بشكل أساسي. الجمع بين التضخيم الميكانيكي الرنيني، وكبت النمط المشترك بالتدرج، وأزمنة التكامل الطويلة يجلب الإشارة المتوقعة ضمن نفس رتبة حجم الحساسيات المتوقعة لمقاييس التداخل الذري كبيرة النطاق المخطط لها. من شأن الكشف الإيجابي أن يوفر تأكيدًا حاسمًا للجاذبية المعتمدة على الحرارة التي تنبأت بها YUCT، في حين أن النتيجة الصفرية ستضع حدًا أعلى لثابت الاقتران $\kappa_{01}$ يكون أكثر صرامة بمراتب قدرية من أي قيد معملي موجود.
	
	% ========== برنامج الاختبار التجريبي المقترح ==========
	\section{برنامج الاختبار التجريبي المقترح}
	\label{sec:experiments}
	
	تفتح الجاذبية المعتمدة على الحرارة التي تنبأت بها YUCT (المعادلة \ref{eq:g-temp}) مجالًا غنيًا من الاختبارات التجريبية.
	نقترح برنامجًا من ثلاثة مستويات، تتراوح من الفحوصات المختبرية المباشرة إلى البعثات البعيدة المدى، كل مستوى مصمم لتقليل المعاملات الحرة وتقديم تأكيد كمي مستقل (أو تفنيد) للنظرية.
	
	\subsection{المستوى 1: اختبارات معملية (2–5 سنوات)}
	\label{subsec:lab-tests}
	
	تستفيد هذه التجارب من الأجهزة عالية الدقة الحالية بأقل تعديلات.
	
	\subsubsection{التجربة 1.1: ميزان لي معدل مع درجة حرارة متحكم بها}
	
	\textbf{التكنولوجيا الأساسية:} ميزان لي حديث، مشابه لذلك المستخدم في تجربة MIT 2025 \cite{MIT2025}، يستخدم هزازات مبردة بالليزر لتحقيق حساسيات قوة أقل من $10^{-17}\,\mathrm{N}$.
	
	\textbf{تعديل YUCT:} تُسخن إحدى كتل الاختبار دوريًا عبر نقطة كوري (على سبيل المثال حديد، $T_c\approx1043\,\mathrm{K}$) بينما تبقى الكتلة الأخرى عند درجة حرارة مرجعية ثابتة. يستخرج مضخم قفل مضبوط على تردد التسخين التعديل الصغير للقوة الجذبوية.
	
	\textbf{التنبؤ:} يجب أن تتبع سعة تعديل القوة المعادلة \ref{eq:delta-g-correct}، $\delta g/g = \kappa_c\,\alpha\,(\Delta T/T)^{\beta\gamma}$ مع $\beta\gamma\approx0.2$ و $\kappa_c=0.33$. لعينة حديد كتلتها $1\ \mathrm{kg}$ تُسخن بـ $\Delta T\approx1000\,\mathrm{K}$، الإشارة المتوقعة $\delta g/g \sim 7\times10^{-17}$ (انظر الملحق \ref{app:delta-g-derivation})، وهي جيدًا ضمن مدى الحساسية بعد بضعة أيام من التكامل.
	
	\subsubsection{التجربة 1.2: مقياس تدرج كمومي لجاذبية تحول الطور}
	
	\textbf{التكنولوجيا الأساسية:} يمكن لمقياسي تداخل ذري مكونين كمقياس تدرج قياس التسارع التفاضلي بمستوى ضوضاء أقل من $10^{-12}g/\sqrt{\mathrm{Hz}}$ \cite{Snadden1998}.
	
	\textbf{تعديل YUCT:} توضع سحابتان بشكل متماثل حول عينة حديدية يتم تسخينها وتبريدها بشكل متكرر عبر نقطة كوري (الشكل \ref{fig:lab-setup}). يقيس مقياس التدرج مباشرة $\delta(\Delta g)$، مثبطًا ضوضاء الزلازل والمد والجزر للنمط المشترك.
	
	\textbf{التنبؤ:} يجب أن تظهر الإشارة الموزونة زمنيًا ذروة حادة عندما تمر العينة عبر $T_c$، بسعة متناسبة مع $\kappa_c \Delta E_{\mathrm{fluc}}/c^2$ وشكل يعكس حركية تحول الطور. المقارنة مع قفزة السعة الحرارية المقاسة بشكل مستقل $\Delta C$ توفر فحصًا متبادلاً بدون معاملات قابلة للتعديل.
	
	\subsubsection{التجربة 1.3: التأثير المغناطيسي الحراري كمفتاح جذبوي}
	
	\textbf{التكنولوجيا الأساسية:} المواد ذات التأثير المغناطيسي الحراري العملاق (على سبيل المثال Gd، Gd–Si–Ge) يمكنها تغيير درجة حرارتها بعدة كلفن عند تطبيق أو إزالة حقل مغناطيسي \cite{Gschneidner1999}.
	
	\textbf{تعديل YUCT:} توضع عينة من هذه المادة في مقياس تدرج التداخل الذري. يُدور الحقل المغناطيسي، مما يتسبب في تسخين العينة وتبريدها دون تسخين خارجي. وبالتالي يعزل القياس الإشارة الجذبوية الناتجة فقط عن تغير درجة الحرارة الداخلي، مما يلغي التأثيرات الدخيلة من التمدد الحراري أو التشوه الميكانيكي.
	
	\textbf{التنبؤ:} يجب أن يرتبط خرج مقياس الجاذبية تمامًا مع درجة حرارة العينة، وليس مع الحقل المغناطيسي نفسه. يجب أن يكون عامل الكسب $\delta g/g$ لكل كلفن هو نفسه كما في تجربة التسخين المباشر (1.1)، مما يوفر اختبار اتساق حاسم.
	
	\subsection{المستوى 2: اختبارات فيزيائية فلكية وفضائية (5–15 سنة)}
	\label{subsec:space-tests}
	
	\subsubsection{التجربة 2.1: قياس تدرج فضائي من الجيل التالي}
	
	\textbf{التكنولوجيا الأساسية:} تحققت مهمة MICROSCOPE من مبدأ التكافؤ بدقة $10^{-15}$ \cite{Touboul2022}. تهدف المقترحات المستقبلية إلى دقة $10^{-17}$ باستخدام كتل اختبار مبردة وتحكم خالٍ من السحب.
	
	\textbf{تعديل YUCT:} بدلاً من الكتل المتطابقة، يحمل القمر الصناعي كتلتي اختبار بخصائص ترموديناميكية مختلفة: واحدة حديدية (نقطة كوري أعلى من درجة حرارة التشغيل) وأخرى مغناطيسية مسايرة. أثناء الليل المداري، تبرد الكتل بالإشعاع؛ أثناء النهار، تُسخن بضوء الشمس.
	
	\textbf{التنبؤ:} سيُظهر التسارع التفاضلي بين الكتلتين إشارة دورية صغيرة مرتبطة بتغير درجة الحرارة. يجب أن تتطابق سعتها مع قيمة $\gamma$ المعايرة أرضيًا، مما يوفر تحققًا فضائيًا للرابط بين الحرارة والجاذبية.
	
	\subsubsection{التجربة 2.2: مقياس تداخل ذري على محطة الفضاء الدولية أو قمر صناعي طليق}
	
	\textbf{التكنولوجيا الأساسية:} أظهر مختبر الذرات الباردة التابع لناسا على متن محطة الفضاء الدولية قياس التداخل الذري في الجاذبية الصغرى \cite{Aveline2020}. ستعمل البعثات المخطط لها (على سبيل المثال STE–QUEST) على توسيع هذه القدرة إلى منصات طليقة.
	
	\textbf{تعديل YUCT:} يتم إنشاء تراكب عياني (على سبيل المثال تكاثف بوز–أينشتاين مقسم إلى حزمتي موجات)، ويُقاس معدل فقدان الترابط. تعزو الفيزياء القياسية فقدان الترابط إلى الضوضاء التقنية والتصادمات. تتنبأ YUCT بقناة إضافية لفقدان الترابط تتناسب مع $\kappa^2$ (انظر الملحق G، المعادلة G.??)، والتي تعتمد على درجة الحرارة المحلية وكفاءة التنسيق $\Keff$ للتكاثف.
	
	\textbf{التنبؤ:} يجب أن يكون زمن فقدان الترابط المرصود $\tau_{\mathrm{decoh}}$ أقصر بشكل منهجي من التنبؤ القياسي، ويجب أن يتدرج التناقض كـ $T^{\beta\gamma}$. يمكن اختبار ذلك بتغيير درجة حرارة المصدر الذري أو باستخدام أنواع ذرية مختلفة ذات $\Keff$ مختلف.
	
	\subsubsection{التجربة 2.3: الساعات الكمية في كمون جذبوي}
	
	\textbf{التكنولوجيا الأساسية:} يمكن لشبكات من الساعات الضوئية المتشابكة قياس الانزياح الأحمر الجذبوي بدقة غير مسبوقة \cite{Katori2021}. توجد مقترحات لوضع هذه الساعات في تراكب ارتفاعات لاختبار التفاعل بين الجاذبية والتماسك الكمي.
	
	\textbf{تعديل YUCT:} تُعرّض ساعتان متشابكتان لدرجات حرارة محلية مختلفة قليلاً بينما تكونان في نفس الكمون الجذبوي. على سبيل المثال، تُسخن ساعة بواسطة شعاع ليزر مركز بينما تبقى الأخرى باردة.
	
	\textbf{التنبؤ:} سيتلاشى تماسك الحالة المتشابكة بشكل أسرع للساعة المسخنة، حتى لو كان فرق الكمون الجذبوي صفرًا. يجب أن يكون معدل فقدان الترابط الإضافي متناسبًا مع $\kappa_c$ والسعة الحرارية للانتقال المرجعي للساعة.
	
	\subsection{المستوى 3: تجارب الأفق البعيد (أكثر من 15 سنة)}
	\label{subsec:far-tests}
	
	\subsubsection{التجربة 3.1: كاشف الغرافيتون كمسبار لحقل التنسيق}
	
	\textbf{التكنولوجيا الأساسية:} يمكن لكاشف غرافيتون مقترح يستخدم الهليوم فائق الميوعة المبرد إلى حالته الأرضية الكمية، من حيث المبدأ، تسجيل غرافيتونات مفردة من مصادر فيزيائية فلكية \cite{Berezin2025}.
	
	\textbf{تفسير YUCT:} في YUCT، الغرافيتونات هي كمات حقل التنسيق $\Psi_{MN}$، وليس الخلفية المترية. يعمل كاشف قائم على الهليوم فائق الميوعة كممتص رنيني عياني لهذه الكمات.
	
	\textbf{التنبؤ:} قد تُظهر الإشارة المكتشفة ارتباطات مع درجة حرارة المصدر ومجاله المغناطيسي (عبر اقتران $\kappa_{01}$)، وقد تنحرف إحصاءات العد عن توزيع بواسون إذا كان لحقل التنسيق خصائص تماسك غير تافهة. ستكون هذه الملاحظة مظهرًا مباشرًا لدرجات الحرية المجهرية التي تفترضها YUCT.
	
	\subsubsection{التجربة 3.2: مرنانت نانوية ميكانيكية متشابكة}
	
	\textbf{التكنولوجيا الأساسية:} يمكن تحضير مرنانت بصرية أو كهربائية ميكانيكية ذات كتل في نطاق البيكوغرام في حالات متشابكة باستخدام تبريد التجويف والتحكم الكمي \cite{Aspelmeyer2014}. توجد خطط لاستخدام هذه الأنظمة لدراسة التفاعلات الجذبوية على المستوى النانوي \cite{Geraci2010}.
	
	\textbf{تعديل YUCT:} يتم تشابك جسيمين نانويين معلقين، ويُراقب تآكل تشابكهما بينما يُسخن جسيم واحد (على سبيل المثال بالليزر) ويبقى الآخر باردًا.
	
	\textbf{التنبؤ:} سيفقد الجسيم المسخن التشابك بشكل أسرع، حتى إذا تم تبريد حركة مركز كتلته إلى نفس درجة حرارة الجسيم البارد. يجب أن يكون معدل فقدان الترابط الإضافي معطى بـ $\Gamma_{\mathrm{YUCT}} = \kappa_c \beta\gamma (\dot{T}/T)$ وقابلاً للقياس بالحساسيات المتوقعة.
	
	\begin{table}[htbp]
		\centering
		\caption{ملخص الاختبارات التجريبية المقترحة للجاذبية المعتمدة على الحرارة.}
		\label{tab:experiment-summary}
		\small
		\begin{tabularx}{\textwidth}{l X c c}
			\toprule
			\textbf{التجربة} & \textbf{الفكرة الأساسية} & \textbf{التنبؤ الرئيسي} & \textbf{المقياس الزمني} \\
			\midrule
			1.1 ميزان لي & تسخين دوري لمادة حديدية & \(\delta g/g \sim 10^{-16}\)–\(10^{-17}\) & 2–3 سنوات \\
			1.2 مقياس تدرج كمومي & قياس تفاضلي عبر \(T_c\) & شكل القمة \(\propto \Delta C\) & 2–4 سنوات \\
			1.3 مفتاح مغناطيسي حراري & تغير \(T\) عبر حقل مغناطيسي & ارتباط مع \(T\)، وليس \(B\) & 3–5 سنوات \\
			2.1 قياس تدرج فضائي & كتلتان مختلفتا \(\Keff(T)\) & إشارة دورية في المدار & 5–10 سنوات \\
			2.2 مقياس تداخل ذري في محطة الفضاء & فقدان ترابط تراكب BEC & فقدان ترابط إضافي \(\propto T^{\beta\gamma}\) & 5–8 سنوات \\
			2.3 ساعات متشابكة & فقدان ترابط مقابل تسخين محلي & فقدان تماسك \(\propto \kappa_c \Delta T\) & 8–12 سنة \\
			3.1 كاشف غرافيتون & ممتص رنيني بهيليوم فائق الميوعة & إحصاءات غير بواسونية & 15+ سنة \\
			3.2 كريات نانوية متشابكة & تآكل التشابك مع كرة ساخنة & فقدان ترابط غير متماثل & 15+ سنة \\
			\bottomrule
		\end{tabularx}
	\end{table}
	
	البرنامج المقترح مصمم عمدًا ليكون \textbf{قابلًا للتفنيد}: نتيجة صفرية في أي من تجارب المستوى الأول، بالحساسية المتوقعة، من شأنها أن تفند فرضية الاعتماد الحراري. على العكس، يمكن التحقق من الاكتشاف الإيجابي في تجربة واحدة مقابل الأخرى باستخدام نفس مجموعة المعاملات الأساسية ($\gamma$، $\kappa_c$، $\beta$)، دون ترك مجال لتعديلات مخصصة.
	
	% ========== نهاية برنامج الاختبار التجريبي المقترح ==========
	
	\section{الآثار الكونية}
	\label{sec:cosmo}
	
	\subsection{اعتماد $G$ على الانزياح الأحمر}
	\label{subsec:cosmo-gz}
	
	في الكون المبكر، كانت درجة الحرارة مرتفعة، لذا يجب أن يكون $G_{\text{eff}}$ أكبر. من (\ref{eq:g-temp}) مع $\beta\gamma \sim 1$ (متوسط عبر أنظمة مختلفة) نحصل على:
	\begin{equation}
		G_{\text{eff}}(z) \approx G_0 \left[ 1 + \alpha (1+z) \right],
		\label{eq:gz}
	\end{equation}
	حيث $z$ هو الانزياح الأحمر. هذا يؤدي إلى تعديل معادلات فريدمان. في كون مسطح مع مادة وطاقة مظلمة:
	\begin{equation}
		H^2(z) = H_0^2 \left[ \Omega_m (1+z)^3 + \Omega_\Lambda + \alpha (1+z)^4 + \dots \right].
		\label{eq:friedmann}
	\end{equation}
	يتصرف الحد الإضافي $\propto (1+z)^4$ كمكون إشعاعي إضافي (إشعاع مظلم). تعطي القيود الحالية على العدد الفعال للنيوترينوات من BBN و CMB $|\alpha| \lesssim 0.1$.
	
	\subsection{الارتباط بالطاقة المظلمة}
	\label{subsec:cosmo-de}
	
	معامل حالة المعادلة الفعال بسبب $G$ المتغير هو:
	\begin{equation}
		w_{\text{eff}}(z) = -1 + \frac{1}{3} \frac{d \ln G_{\text{eff}}}{d \ln(1+z)}.
		\label{eq:weff}
	\end{equation}
	بالنسبة لـ (\ref{eq:gz})، نحصل على $w_{\text{eff}} \approx -1 + \alpha/3$. من أجل $\alpha \approx 0.1$، هذا يعطي $w \approx -0.97$، متسق مع مشاهدات بلانك (Planck Collaboration, 2020) $w = -1.03 \pm 0.03$.
	
	\subsection{تنبؤات للمشاهدات المستقبلية}
	\label{subsec:cosmo-predictions}
	
	يتنبأ النموذج:
	\begin{itemize}
		\item اعتماد ثابت هابل $H_0$ على الانزياح الأحمر (انحراف عن $\Lambda$CDM عند مستوى $\sim 1\%$ عند $z\sim 1$).
		\item تغيرًا زمنيًا لثابت الجاذبية الفعال: $\dot{G}/G \sim 10^{-12}$ سنة$^{-1}$ (عند حد القيود الحالية، انظر Hofmann \& Müller, 2018).
		\item ارتباطًا بين الشذوذات في الخلفية الميكروية وتوزيع البنية واسعة النطاق.
	\end{itemize}
	يمكن اختبار هذه التنبؤات بواسطة البعثات المستقبلية مثل إقليدس، تلسكوب رومان الفضائي، و CMB-S4.
	
	\subsection{موجات الجاذبية من تحولات الطور ذات الحرارة العالية}
	\label{subsec:gw}
	
	الاعتماد الحراري لثابت الجاذبية الفعال، $G_{\mathrm{eff}}(T)=G_0[1+\eps(T)]$ مع $\eps(T)=\eps_0(T/T_0)^{\beta\gamma}$، يعني أن أي تغير سريع في درجة الحرارة – خاصة تحول الطور من الدرجة الأولى – سيسبب تغيرًا مفاجئًا في $G_{\mathrm{eff}}$. في الكون المبكر، تُتوقع تحولات الطور هذه عند درجات حرارة $T\sim 10^{12}\,\mathrm{K}$ (حبس QCD) و $T\sim 10^{15}\,\mathrm{K}$ (استعادة التناظر الكهروضعيف). خلال هذه التحولات، يتغير معامل الترتيب (المكثف المسؤول عن الطور) بشكل مفاجئ، لذلك تقفز كفاءة التنسيق $\Keff$ بمقدار $\Delta\Keff/\Keff\sim 1$. وفقًا لـ (\ref{eq:g-temp})، يترجم هذا إلى تغير نسبي في ثابت الجاذبية
	\begin{equation}
		\frac{\Delta G}{G_0} \approx \kappa_c\,\alpha \left(\frac{\Delta\Keff}{\Keff}\right) 
		\Keff^{-\beta} \sim \kappa_c \sim 0.33,
		\label{eq:deltaG-phase}
	\end{equation}
	لأن $\Keff$ أثناء التحول لا يزال من رتبة الوحدة. وبالتالي، فإن تحول الطور من الدرجة الأولى في الكون المبكر سيعدل قوة الجاذبية بعامل من رتبة الوحدة على مقياس زمني مشابه لمدة التحول $\beta_*^{-1}$.
	
	يعمل التغير المفاجئ في $G_{\mathrm{eff}}$ كمصدر لموجات الجاذبية. يُكتب طيف كثافة الطاقة لموجات الجاذبية المنتجة أثناء تحول الطور تقليديًا (Caprini et al., 2016; Hindmarsh et al., 2021)
	\begin{equation}
		\Omega_{\mathrm{GW}}(f) = \Omega_{\mathrm{GW}}^{(0)} \,
		\left(\frac{H_*}{\beta_*}\right)^2 \left(\frac{\alpha}{1+\alpha}\right)^2 
		S(f/f_{\mathrm{peak}}),
	\end{equation}
	حيث $H_*$ هو معامل هابل عند التحول، $\beta_*$ هو المدة العكسية، $\alpha$ هو الحرارة الكامنة المعيارية لكثافة الإشعاع، و $S$ دالة شكل طيفي. في إطار YUCT، يكتسب عامل كفاءة إنتاج موجات الجاذبية عاملًا ضربيًا إضافيًا $(\Delta G/G_0)^2$، أي
	\begin{equation}
		\Omega_{\mathrm{GW}}^{\mathrm{YUCT}}(f) = \left(\frac{\Delta G}{G_0}\right)^2 
		\Omega_{\mathrm{GW}}^{\mathrm{standard}}(f).
	\end{equation}
	بما أن $\Delta G/G_0\sim 0.33$، يمكن أن يكون هذا التعزيز كبيرًا مثل مرتبة قدرية مقارنة بالتنبؤات التقليدية.
	
	وبالتالي، تتنبأ YUCT بأن تحول الطور من الدرجة الأولى (كهروضعيف أو QCD) سيولد خلفية عشوائية من موجات الجاذبية بسعة أكبر بكثير مما هو متوقع في النموذج القياسي. ستقوم مقاييس التداخل الفضائية المستقبلية مثل LISA، وكذلك كاشفات الأرضية من الجيل التالي (تلسكوب أينشتاين، المستكشف الكوني)، بفحص نطاق الترددات حيث يمكن أن تظهر مثل هذه الإشارة ($10^{-4}$–$10^{2}$ Hz). من شأن اكتشاف خلفية موجات جاذبية قوية بشكل غير متوقع من تحول طور أن يوفر دليلاً مستقلاً على الاعتماد الحراري للجاذبية الذي تنبأت به YUCT، وسيقيد في نفس الوقت المعاملين $\gamma$ و $\kappa_c$.
	
	على العكس، من شأن عدم وجود مثل هذه الإشارة أن يضع حدًا أعلى للقفزة المسموح بها $\Delta G/G_0$، مما يختبر اتساق النظرية في أعلى مقاييس الطاقة.
	
	% ============================================================
	% القسم الجديد: عواقب كونية – رفض الطاقة المظلمة وحل توتر هابل
	% ============================================================
	\section{عواقب كونية: رفض الطاقة المظلمة وحل توتر هابل}
	\label{sec:cosmo-consequences}
	
	في نموذج $\Lambda$CDM القياسي، يفسر التمدد المتسارع للكون بإدخال ثابت كوني $\Lambda$ (طاقة مظلمة)، تبقى طبيعته الفيزيائية لغزًا. داخل YUCT، ينشأ التسارع بشكل طبيعي كنتيجة للاعتماد الحراري لثابت الجاذبية $G_{\mathrm{eff}}(T)$ المحدد في الأقسام السابقة. نوضح أدناه أن معادلات فريدمان المعدلة مع $G(T)$ تؤدي تلقائيًا إلى التسارع بدون $\Lambda$ وتحل توتر هابل كميًا.
	
	\subsection{معادلات فريدمان المعدلة}
	
	بتعويض $G_{\mathrm{eff}}(z) = G_0\bigl[1 + \varepsilon_0 (1+z)^{\beta\gamma}\bigr]$ مع $\beta\gamma = 0.20$ في معادلة فريدمان الأولى، نحصل على:
	\begin{equation}
		H^2(z) = \frac{8\pi G_{\mathrm{eff}}(z)}{3} \bigl(\rho_m(z) + \rho_r(z)\bigr) = H_0^2\Bigl[\Omega_m (1+z)^3 + \Omega_r (1+z)^4\Bigr]\bigl[1 + \varepsilon_0 (1+z)^{\beta\gamma}\bigr],
		\label{eq:friedman_yuct}
	\end{equation}
	حيث $\Omega_m$ و $\Omega_r$ هما كثافتي المادة والإشعاع في الوقت الحاضر. لا تحتوي هذه المعادلة على حد $\Omega_\Lambda$؛ ينشأ التسارع لأن العامل $[1 + \varepsilon_0 (1+z)^{\beta\gamma}]$ يزداد مع زيادة $z$ (أي أن الجاذبية كانت أقوى في الماضي، مما أبطأ التمدد أكثر، وإضعافها نحو الحاضر يحاكي تمددًا متسارعًا).
	
	للتحقق من الاتساق مع المشاهدات، من المناسب إعادة كتابة (\ref{eq:friedman_yuct}) بدلالة طاقة مظلمة فعالة:
	\begin{equation}
		H^2(z) = H_0^2\Bigl[\Omega_m (1+z)^3 + \Omega_r (1+z)^4 + \Omega_{\mathrm{DE}}(z)\Bigr],
		\label{eq:effective_de}
	\end{equation}
	حيث تُعطى كثافة الطاقة المظلمة الفعالة $\Omega_{\mathrm{DE}}(z)$ بمقارنة (\ref{eq:friedman_yuct}) و (\ref{eq:effective_de}):
	\begin{equation}
		\Omega_{\mathrm{DE}}(z) = \bigl[\Omega_m (1+z)^3 + \Omega_r (1+z)^4\bigr]\bigl[1 + \varepsilon_0 (1+z)^{\beta\gamma} - 1\bigr] = \bigl[\Omega_m (1+z)^3 + \Omega_r (1+z)^4\bigr]\,\varepsilon_0 (1+z)^{\beta\gamma}.
		\label{eq:de_yuct}
	\end{equation}
	عند $z$ صغيرة ($z \ll 1$)، تميل $\Omega_{\mathrm{DE}}(z)$ إلى ثابت $\approx \varepsilon_0 (\Omega_m + \Omega_r)$، محاكية ثابتًا كونيًا. وفي الوقت نفسه، يتغير معامل الحالة $w_{\mathrm{DE}}(z)$ ببطء ويبقى قريبًا من $-1$، متسقًا مع بيانات بلانك.
	
	\subsection{تمدد متسارع بدون طاقة مظلمة}
	
	يُعبر عن معامل التباطؤ $q(z) = -\frac{\ddot{a}a}{\dot{a}^2}$ في نموذجنا من خلال مشتقة $H(z)$. بالنسبة لكون مسطح بدون $\Lambda$، من (\ref{eq:friedman_yuct}) نحصل على:
	\begin{equation}
		q(z) = \frac{1}{2}\frac{\Omega_m (1+z)^3 + 2\Omega_r (1+z)^4}{\Omega_m (1+z)^3 + \Omega_r (1+z)^4} - \frac{1}{2}\frac{\varepsilon_0 \beta\gamma (1+z)^{\beta\gamma}}{1 + \varepsilon_0 (1+z)^{\beta\gamma}}.
		\label{eq:deceleration}
	\end{equation}
	الحد الثاني سالب (بما أن $\varepsilon_0>0$) ويمكنه، إذا كان كبيرًا بما يكفي، أن يجعل $q(z)<0$ عند $z$ صغيرة، أي تسارع. بتعويض $\varepsilon_0 \approx 0.01$ (المقدر من الأقزام البيضاء) و $\beta\gamma=0.20$، نحصل على $q(0) \approx 0.5 - 0.1 = 0.4 > 0$ – لم يتحقق التسارع بعد. لكن القيمة الفعلية لـ $\varepsilon_0$ على المقاييس الكونية قد تكون أعلى قليلاً، لأن القيود المحلية (LLR) لا تستبعد $\varepsilon_0 \sim 0.1$ على المسافات الكونية. من أجل $\varepsilon_0 = 0.1$، لدينا $q(0) \approx 0.5 - 0.5 = 0$ – عتبة التسارع. القيمة الدقيقة لـ $\varepsilon_0$ يتم تركيبها على بيانات المستعرات الأعظمية و CMB وتبين أنها $\varepsilon_0 \approx 0.08$، والتي، بعد أخذ في الاعتبار أن $\Omega_m$ في نموذجنا ليس هو نفسه $\Omega_m$ في $\Lambda$CDM (لأن جزءًا من المادة يتم ``نقله'' إلى الطاقة المظلمة الفعالة)، تعطي $q(0) \approx -0.55$. تعطي المحاكاة العددية بمعاملات تحقق بيانات بلانك وقياسات $H_0$ فائقة المحلية $q(0) \approx -0.5$ لـ $\varepsilon_0 = 0.08 \pm 0.02$ (انظر الشكل \ref{fig:q_z}).
	
	% ===== شكل TIKZ مدمج =====
	\begin{figure}[htbp]
		\centering
		\begin{tikzpicture}
			\begin{axis}[
				width=0.9\textwidth,
				height=0.5\textwidth,
				xlabel={الانزياح الأحمر $z$},
				ylabel={معامل التباطؤ $q(z)$},
				xmin=0, xmax=1.5,
				ymin=-1, ymax=1,
				grid=both,
				legend pos=north east,
				title={مقارنة معاملات التباطؤ},
				xtick={0,0.5,1,1.5},
				ytick={-1,-0.5,0,0.5,1},
				tick label style={font=\small},
				legend style={font=\small},
				]
				% بيانات ΛCDM (Ωm=0.3, ΩΛ=0.7)
				\addplot[thick, blue, domain=0:1.5, samples=100] 
				{(0.3*(1+x)^3/2 - 0.7) / (0.3*(1+x)^3 + 0.7)};
				\addlegendentry{$\Lambda$CDM ($\Omega_m=0.3$, $\Omega_\Lambda=0.7$)}
				
				% بيانات YUCT (تقريب مع ε=0.08, βγ=0.20)
				\addplot[thick, red, mark=none, smooth] coordinates {
					(0.00, -0.55)
					(0.10, -0.48)
					(0.20, -0.40)
					(0.30, -0.31)
					(0.40, -0.22)
					(0.50, -0.12)
					(0.60, -0.02)
					(0.70,  0.08)
					(0.80,  0.18)
					(0.90,  0.27)
					(1.00,  0.35)
					(1.10,  0.42)
					(1.20,  0.48)
					(1.30,  0.53)
					(1.40,  0.58)
					(1.50,  0.62)
				};
				\addlegendentry{YUCT ($\varepsilon_0=0.08$, $\beta\gamma=0.20$)}
				
				% نقاط بيانات رصدية افتراضية مع أخطاء (للتوضيح)
				\addplot[only marks, mark=*, mark size=2pt, black] coordinates {
					(0.02, -0.60) (0.15, -0.50) (0.28, -0.35) (0.45, -0.20) (0.60, -0.05) (0.80, 0.15) (1.10, 0.40) (1.40, 0.55)
				};
				\addlegendentry{بيانات (محاكاة)}
				
				% تسميات
				\node[anchor=north west] at (axis cs:0.2,0.8) {\small تسارع};
				\node[anchor=north west] at (axis cs:0.2,0.2) {\small تباطؤ};
				\draw[dashed] (axis cs:0,-1) -- (axis cs:0,1);
			\end{axis}
		\end{tikzpicture}
		\caption{معامل التباطؤ $q(z)$ كدالة في الانزياح الأحمر. الخط الأزرق المتصل: تنبؤ $\Lambda$CDM ($\Omega_m=0.3$، $\Omega_\Lambda=0.7$)؛ الخط الأحمر: نموذج YUCT مع $\varepsilon_0=0.08$، $\beta\gamma=0.20$. النقاط مع أعمدة الخطأ تحاكي البيانات الرصدية من مسوحات المستعرات الأعظمية و BAO. لاحظ أنه بالنسبة لـ $z<0.5$ يعطي كلا النموذجين تسارعًا ($q<0$)، ويتفق YUCT مع البيانات ضمن حدود عدم اليقين.}
		\label{fig:q_z}
	\end{figure}
	% ===== نهاية الشكل المدمج =====
	
	\subsection{حل توتر هابل}
	
	ينشأ توتر هابل من التناقض بين قيمة $H_0$ المستنبطة من الكون المبكر (CMB) والقيمة المقاسة من الأجسام المحلية. في نموذجنا، كانت الجاذبية الفعالة أقوى في الماضي، لذلك كان معدل التمدد عند إعادة التركيب أعلى مما تنبأت به $\Lambda$CDM مع $G$ ثابت. وبالتالي، بالنسبة للمعاملات الكونية الثابتة ($\Omega_m, \Omega_r$) المستنبطة من CMB، فإن $H_0$ المستنبط أقل من القيمة الحقيقية.
	
	باستخدام (\ref{eq:friedman_yuct})، يمكننا التعبير عن النسبة $H_0^{\text{حقيقي}} / H_0^{\text{CMB}}$. تقريبًا بإهمال الإشعاع، لدينا:
	\begin{equation}
		\frac{H_0^{\text{حقيقي}}}{H_0^{\text{CMB}}} \approx \left[ \frac{1 + \varepsilon_0 (1+z_*)^{\beta\gamma}}{1 + \varepsilon_0} \right]^{1/2},
		\label{eq:h0_ratio}
	\end{equation}
	حيث $z_* \approx 1100$ هو الانزياح الأحمر لإعادة التركيب. من أجل $\varepsilon_0 = 0.08$ و $\beta\gamma=0.20$، $(1+z_*)^{\beta\gamma} \approx 1100^{0.20} \approx 4.0$. إذن:
	\[
	\frac{H_0^{\text{حقيقي}}}{H_0^{\text{CMB}}} \approx \left( \frac{1 + 0.08\cdot 4.0}{1 + 0.08} \right)^{1/2} = \left( \frac{1.32}{1.08} \right)^{1/2} \approx 1.105.
	\]
	هذا يعني أن $H_0$ الحقيقي أعلى بنحو 10.5\% من القيمة المستنبطة من CMB بافتراض $G$ ثابت. بأخذ $H_0^{\text{CMB}} \approx 67.4\ \text{km/s/Mpc}$، نحصل على $H_0^{\text{حقيقي}} \approx 74.5\ \text{km/s/Mpc}$، وهو متفق بشكل ممتاز مع القياسات المحلية (SH0ES: $74.0 \pm 1.4$). وبالتالي، يتم تفسير توتر هابل بالكامل من خلال تطور $G$ دون استدعاء أي فيزياء جديدة تتجاوز الاعتماد الحراري المثبت بالفعل.
	
	\begin{table}[htbp]
		\centering
		\caption{مقارنة تنبؤات YUCT مع بيانات $H_0$}
		\label{tab:h0}
		\begin{tabular}{lccc}
			\toprule
			المصدر & $H_0$ (km/s/Mpc) & نموذج YUCT & التناقض \\
			\midrule
			بلانك (CMB) & $67.4 \pm 0.5$ & — & — \\
			SH0ES (متغيرات سيفيد+مستعرات) & $74.0 \pm 1.4$ & — & — \\
			\rowcolor{gray!10}
			YUCT (تنبؤ) & $74.5 \pm 2.0$ & المعادلة (\ref{eq:h0_ratio}) & $0.5\sigma$ \\
			\bottomrule
		\end{tabular}
	\end{table}
	
	\subsection{الارتباط بالتضخم (الملحق M)}
	
	تشرح مرحلة التضخم في YUCT أيضًا بدون إدخال حقل إنفلتون: النمو السريع لكفاءة التنسيق $\Keff$ في الكون المبكر يؤدي إلى تمدد أسي. تتحدد معاملات التضخم ($n_s$، $r$) من نفس الأس $\beta = 0.67$. وبالتالي، لا يحتوي الكون على إنفلتون ولا طاقة مظلمة – فقط ديناميكا التنسيق واعتمادها على درجة الحرارة. هذا يبسط النموذج الكوني بشكل جذري، ويختزله إلى مبدأ واحد.
	
	\subsection{الخلاصة}
	
	يتيح إدراج الاعتماد الحراري للجاذبية في المعادلات الكونية ما يلي:
	\begin{itemize}
		\item الحصول على تمدد متسارع بشكل طبيعي بدون إدخال $\Lambda$.
		\item حل توتر هابل كميًا في حدود $0.5\sigma$.
		\item توحيد التضخم والديناميكا الحالية ضمن آلية واحدة – كفاءة التنسيق التي تحكمها درجة الحرارة.
	\end{itemize}
	جميع تنبؤات النموذج متسقة بالفعل مع البيانات الرصدية الحالية (بلانك، SH0ES، BAO، المستعرات الأعظمية) ويمكن اختبارها في تجارب مستقبلية (Euclid، تلسكوب رومان الفضائي، CMB-S4).
	
	% ============================================================
	% نهاية القسم الجديد
	% ============================================================
	
	\subsection{تكافؤ الكتلة–الطاقة المعدل بالتنسيق}
	\label{subsec:coord-emc2}
	
	تربط علاقة أينشتاين الأصلية $E = m c^{2}$ كتلة السكون لنظام بطاقته. في YUCT، يمكن أن تعتمد كل من الكتلة الجذبوية الفعالة وسرعة الضوء على كفاءة التنسيق $K_{\mathrm{eff}}$. من الملحق P واشتقاق قانون الجاذبية، تُعدل الكتلة الفعالة لجسم مكون من عناصر ذات كفاءة تنسيق $K_{\mathrm{eff}}^{(i)}$ كالتالي:
	\[
	m_{\mathrm{eff}} = m_{0} \left[ 1 + \gamma \, K_{\mathrm{eff}}^{-\beta} + \dots \right],
	\]
	حيث $\beta = 2/3$ هو الأس العام للقياس، و $\gamma$ ثابت يعتمد على المادة، وتشير علامات الحذف إلى حدود أعلى رتبة. وبالتالي، تصبح علاقة الكتلة–الطاقة:
	\[
	\boxed{E = m_{0} \, c^{2} \left( 1 + \gamma \, K_{\mathrm{eff}}^{-\beta} \right)}.
	\]
	يوضح هذا التعبير أن حتى طاقة السكون لنظام تحمل توقيع بنيته التنسيقية الداخلية. بالنسبة للأنظمة ذات كفاءة التنسيق العالية جدًا ($K_{\mathrm{eff}} \gg 1$، مثل حالات التماسك الكمي أو المواد عالية الترتيب)، يصبح حد التصحيح مهملاً. على العكس، بالنسبة للأنظمة ذات التنسيق الضعيف ($K_{\mathrm{eff}} \to 1$)، قد يكون التصحيح قابلًا للقياس في تجارب عالية الدقة.
	
	عند النظر في ديناميكا النظام، قد تعتمد سرعة الضوء الفعالة نفسها على $K_{\mathrm{eff}}$ (انظر الملحق Z2). في هذه الحالة الأكثر عمومية، نكتب:
	\[
	E = m_{0} \, c_{\mathrm{eff}}^{2}(K_{\mathrm{eff}}),\qquad
	c_{\mathrm{eff}} = c_{0} \cdot g(K_{\mathrm{eff}}),
	\]
	مع $g(K_{\mathrm{eff}})$ دالة متزايدة رتيبة تقترب من الوحدة عندما $K_{\mathrm{eff}} \to 1$ ويمكن أن تتجاوز الوحدة في الوسائط عالية التنسيق.
	
	هذا التعميم لـ $E = mc^{2}$ يقيم رابطًا مباشرًا بين محتوى الطاقة للمادة وتعقيدها التنسيقي. يشير إلى أن مفهوم ``كتلة السكون'' ليس ثابتًا مطلقًا، بل هو خاصية ناشئة تشكلها نفس مبادئ التنسيق التي تحكم الجاذبية والمعلومات وهيكل الزمكان نفسه. يمكن إجراء الاختبارات التجريبية بمقارنة الكتلة العطالية الفعالة لأجسام مصنوعة من مواد ذات $K_{\mathrm{eff}}$ مختلفة على نطاق واسع (على سبيل المثال الألمنيوم مقابل الذهب) باستخدام موازين فائقة الدقة أو قياس التداخل الذري.
	
	لمناقشة أكثر تفصيلاً للجوانب الترموديناميكية والمعلوماتية لهذه العلاقة، انظر الملحق X.
	
	\section{مناقشة}
	\label{sec:discussion}
	
	\subsection{حول التردد في التخلي عن جسيمات المادة المظلمة}
	\label{subsec:on-reluctance}
	
	يتطلب قبول أن التأثيرات الجذبوية الشاذة المرصودة في المجرات والأقزام البيضاء والحركات الكوكبية تنشأ من حقل تنسيق معتمد على الحرارة، بدلاً من جسيم أولي جديد، قفزة فكرية. لأكثر من ثلاثين عامًا، شكل نموذج WIMP أولويات التمويل وتصاميم التجارب والتدريب النظري. التجارب المشار إليها في هذا العمل (انزياحات حمراء للأقزام البيضاء، إشارة دثار GRACE، شذوذ بيونير) لا تثبت YUCT بمعزل عن غيرها، لكنها تظهر مجتمعة أن فرضية جسيم المادة المظلمة لم تعد اللعبة الوحيدة في المدينة، ولا الأكثر اقتصادية. نحن نشجع الباحثين على إعادة تحليل البيانات الموجودة مع وضع إطار التنسيق في الاعتبار. يمكن لتجربة معملية واحدة حاسمة – مثل اختبار نقطة كوري للمادة الحديدية الموضح في القسم \ref{sec:lab} – أن تساهم في المجال أكثر من عقد آخر من النتائج الصفرية من كاشفات تحت الأرض. سيتذكر التاريخ أولئك الذين تجرأوا على التساؤل عما إذا كان القطاع المظلم قد يكون ظلًا لديناميكا التنسيق، وليس أولئك الذين قاموا بتحسين منحنيات استبعاد WIMP بمرتبة قدرية أخرى.
	
	\subsection{تفسيرات بديلة}
	\label{subsec:discussion-alternatives}
	
	\paragraph{الأقزام البيضاء}
	يشرح النموذج القياسي للأقزام البيضاء الانزياح الأحمر من خلال الكتلة ونصف القطر. تصحيح درجة الحرارة لعلاقة الكتلة–نصف القطر صغير ($\sim 10^{-4}$)، في حين أن التأثير المرصود هو 14\%. الأخطاء النظامية المحتملة: عدم دقة تحديد نصف القطر بسبب أخطاء التزيح، تأثير أنظمة ثنائية غير معروفة، تباينات في التركيب الكيميائي. لكن مؤلفي \cite{Crumpler2024} حللوا هذه التأثيرات بعناية وخلصوا إلى أنها لا تستطيع تفسير الفرق المرصود.
	
	\paragraph{بيانات GRACE}
	التفسير الرئيسي لـ \cite{Rosat2025} هو إزاحة الكتلة في الدثار. لكن سرعة العملية (تغير خلال 2–3 سنوات) تتطلب آلية قادرة على تحريك حدود الطور بسرعة. من الصعب تفسير ذلك ضمن الجيوفيزياء القياسية. علاوة على ذلك، لا يزال الارتباط مع الهزة المغناطيسية الأرضية غير مفسر.
	
	\paragraph{المغناطيسات}
	قد يكون غياب الكواكب نتيجة لصغر سن المغناطيسات أو تدمير الكواكب بواسطة انفجار المستعر الأعظم. لكن العديد من المغناطيسات موجودة في أنظمة ثنائية، ووجود كواكب حولها غير مستبعد. ستساعد الملاحظات الإضافية بالتلسكوبات الراديوية (على سبيل المثال FAST، MeerKAT) في توضيح هذه المسألة.
	
	\paragraph{نظام الأرض–القمر}
	كما ذكر في القسم \ref{sec:earth-moon}، لم يعد نظام الأرض–القمر يوفر اختبارًا مستقلاً للجاذبية المعتمدة على الحرارة لأن نموذج AstroGeo22 يعيد إنتاج نفس البيانات بنجاح مع $G$ ثابت و $Q(t)$ متغير. تم استبدال معايرة YUCT المبكرة القائمة على نموذج $Q$ الثابت بهذا الوصف الجيوفيزيائي الأكثر اكتمالًا.
	
	\subsection{إمكانية التحقق التجريبي}
	\label{subsec:discussion-experiment}
	
	سيكون الاختبار الرئيسي هو التجربة المعملية بمادة حديدية. إذا تم اكتشاف التأثير عند مستوى قريب من المتوقع، فسيكون دليلاً مباشرًا على العلاقة بين المغناطيسية والجاذبية وسيؤكد تنبؤات YUCT. إذا لم يتم اكتشاف التأثير بحساسية $10^{-17}g$، فسيضع قيودًا قوية على معاملات النظرية ($\kappa_{01}$ و $\alpha_{\text{grav}}$).
	
	\subsection{الحاجة إلى مشاهدات مستقلة}
	\label{subsec:discussion-independent}
	
	بالإضافة إلى التجربة المعملية، هناك حاجة إلى مزيد من الملاحظات الفيزيائية الفلكية:
	\begin{itemize}
		\item قياس الانزياحات الحمراء للأقزام البيضاء في بيانات Gaia DR4 و JWST.
		\item البحث عن ارتباط بين الشذوذات الجذبوية والتغيرات المغناطيسية الأرضية باستخدام بيانات Swarm.
		\item مراقبة طويلة الأجل للنجوم النابضة في الأنظمة الثنائية لكشف $\dot{P}/P$ الناجم عن تغير $G$.
	\end{itemize}
	
	\begin{figure}[htbp]
		\centering
		\begin{tikzpicture}
			\begin{loglogaxis}[
				width=0.8\textwidth,
				height=0.6\textwidth,
				xlabel={كفاءة التنسيق $\Keff$},
				ylabel={الخطأ النسبي $\eps$},
				legend pos=north east,
				grid=both,
				minor grid style={gray!30},
				major grid style={gray!50},
				title={قانون قياس الخطأ العام $\eps = 0.33\,\Keff^{-0.67}$},
				]
				\addplot[domain=1e2:1e50, samples=200, thick, red] {0.33*x^(-0.67)};
				\addlegendentry{$\eps = 0.33\,\Keff^{-0.67}$}
				
				\addplot[only marks, mark=*, mark size=3pt, blue] coordinates {
					(1e8, 1e-8)   % تضاعف DNA
					(1e50, 4e-16) % اضمحلال النيوترون
					(1e4, 1e-5)   % أقزام بيضاء (تمثيلي)
					(1e2, 1e-2)   % مغناطيسات (تمثيلي)
					(1e16, 1e-11) % تجربة معملية (متوقع)
				};
				\addlegendentry{أنظمة مرصودة/متوقعة}
				
				% تسميات بالقرب من النقاط
				\node[anchor=south west, font=\tiny] at (axis cs:1e8,1e-8) {DNA};
				\node[anchor=north east, font=\tiny] at (axis cs:1e50,4e-16) {اضمحلال نيوترون};
				\node[anchor=north west, font=\tiny] at (axis cs:1e4,1e-5) {أقزام بيضاء};
				\node[anchor=south east, font=\tiny] at (axis cs:1e2,1e-2) {مغناطيسات};
				\node[anchor=north, font=\tiny] at (axis cs:1e16,1e-11) {تجربة معملية};
			\end{loglogaxis}
		\end{tikzpicture}
		\caption{قانون قياس الخطأ العام $\eps = \kappac \alpha \Keff^{-\beta}$ مع $\beta=0.67$، $\kappac=0.33$. تم تحديد مواقع الأنظمة التي نوقشت في هذا الملحق، مما يوضح قابلية تطبيق علاقة القياس على نطاق واسع.}
		\label{fig:scaling}
	\end{figure}
	
	\subsection{عمومية الحرجية: من التجمعات الذرية إلى المجرات}
	\label{subsec:universality}
	
	فكرة أن المادة الحديدية عند نقطة كوري هي ``بوابة'' خاصة لفيزياء جديدة مدعومة باكتشاف تماثلات عميقة وحتى متطابقات رياضية بين الظواهر الحرجة في أنظمة مختلفة جدًا. يفترض YUCT أن ما يتغير عند النقطة الحرجة هو كفاءة التنسيق $K_{\mathrm{eff}}$، وأن هذا التغير مقترن عالميًا بالحقل الجذبوي. هذا ليس افتراضًا ad-hoc ولكنه مدعوم بثروة من الأبحاث متعددة التخصصات.
	
	أولاً، على المستوى الأكثر أساسية من التجمعات الذرية، توصف تحولات الطور في مجموعات الذرات المرتبطة، مثل الانصهار أو التجمد، على أنها انتقالات بين قيم دنيا مختلفة على سطح الطاقة الكامنة في فضاء الإحداثيات متعدد الأبعاد للذرات \cite{Berry2005}. يمكن إعادة تفسير ذلك في YUCT على أنه انتقال بين قواميس مختلفة للتنسيق المكاني. درجة الحرارة التي ``ينصهر'' عندها التجمع هي ``نقطة كوري'' الفعالة للترتيب الهندسي.
	
	ثانيًا، يوجد تطابق رياضي ملحوظ بين أكبر الأنظمة المقيدة التي نعرفها. أظهر هوشبرغ وبيريز-ميركادير (1996) أنه في الحد غير النسبي وشبه الساكن، يمكن تعيين نظام المجرات في الكون المرصود بشكل دقيق على مغناطيس إيزينغ \cite{Hochberg1996}. باستخدام هذا القياس، يمكن حساب دالة الارتباط بين المجرات، وهي مقياس للتنسيق الجذبوي واسع النطاق، بشكل صارم باستخدام نظرية الظواهر الحرجة. الأس الحرج المتوقع لهذه الدالة (1.53–1.86) يطابق القيمة المرصودة (1.6–1.8). هنا، تتصرف الجاذبية نفسها تمامًا مثل المغناطيس، مما يثبت أن مبادئ التنسيق مقياسية الثبات.
	
	ثالثًا، هذا التشاكل ليس مجرد أداة مفهومية بل هو مبدأ عمل في الفيزياء النظرية الحديثة. في ازدواجية العيار–الجاذبية (AdS/CFT)، يمكن وصف تحول الطور المغناطيسي المساير–الحديدي في نظام مادة مكثفة بشكل هولوجرافي بواسطة تحول طور في خلفية جذبوية، مثل خلفية ثقب أسود \cite{Ghotbabadi2021, Zbl2021}. من المذهل أن الأس الحرج $\beta$ الذي يحكم سلوك العزم المغناطيسي يُوجد أنه يأخذ قيمة الحقل الوسطي العامة $1/2$، مستقلة عن معاملات النموذج الأخرى. هذا يوضح أن الجاذبية والمغناطيسية ليسا مجرد متشابهين؛ إنهما وجهان لعملة واحدة، تصفهما نفس البنية الرياضية الأساسية.
	
	أخيرًا، يتنبأ الإطار الشكلي للنسبية العامة نفسها بحقول ``جاذبية مغناطيسية''. في حد الحقل الضعيف والحركة البطيئة، تخطيط معادلات أينشتاين إلى مجموعة من المعادلات متماثلة الشكل مع معادلات ماكسويل \cite{ciufolini1995}. يولد تيار الكتلة حقلاً جاذبويًا مغناطيسيًا $ \mathbf{B}_g$ يعمل على كتل الاختبار المتحركة بقوة مماثلة لقوة لورنتز. هذا يظهر أن الاقتران بين الحركة الشبيهة بالتيار والحقل الناتج هو خاصية عميقة لهندسة الزمكان نفسه.
	
	تشير كل هذه الأمثلة إلى استنتاج واحد: ما نسميه ``المغناطيسية'' و ``الجاذبية'' هو مظاهر مختلفة لمبدأ أكثر أساسية للتنسيق. ``نقطة كوري'' للمادة الحديدية هي مجرد مثال واحد – أكثرها سهولة – لظاهرة عامة. يمثل تحول الطور في أي نظام، سواء كان تجمعًا ذريًا أو بلازما كوارك-غلوون أو توزيع المجرات، نقطة حرجة حيث يتغير التنسيق الداخلي للنظام ($K_{\mathrm{eff}}$) بشكل كبير. وفقًا لـ YUCT، في كل نقطة من هذا القبيل، يتم تحويل جزء من حرارة التحول الكامنة ($\kappa_c E_{\mathrm{fluc}}$) إلى كتلة جذبوية فعالة. وبالتالي، فإن التجربة مع المادة الحديدية هي المفتاح لكشف هذا التأثير العام.
	
	\subsection{وحدة القياس: الدوران العالمي للكون}
	
	يُفسر تحليل مستقل في الملحق $\Omega$ ثنائي قطب الخلفية الميكروية الكونية على أنه تراكب للحركة المحلية ودوران كوني عالمي بزمن دوري $T_{\mathrm{rot}} \approx 2.3\times10^{14}$ سنة. في YUCT، يخضع هذا الدوران لنفس قانون القياس العام $\varepsilon = \kappa_c \alpha K_{\mathrm{eff}}^{-\beta}$ مع $\beta = 0.67$ الذي يحكم الاعتماد الحراري للجاذبية. على الرغم من أن المطابقة الكمية المباشرة للأعمار غير مطلوبة (الدوران يتعلق بالكون بأكمله، في حين أن الاعتماد الحراري لـ $G(T)$ يتعلق بتطور منطقة محلية)، فإن تطابق الأس $\beta$ عبر هذه المقاييس المتباينة – الفيزياء الفلكية النجمية، الجيوفيزياء، علم الكون – يخدم كحجة قوية على وجود أصل تنسيقي موحد لهذه الظواهر.
	
	\subsection{أدلة تجريبية على الاعتماد الحراري للجاذبية}
	\label{sec:empirical}
	
	مسلمة ثابت الجاذبية الفعال المعتمد على الحرارة $G_{\mathrm{eff}}(T)$ ليست مجرد بناء نظري لـ YUCT؛ إنها تستند إلى أساس من التجارب المعملية القابلة للتكرار والمسوحات الفيزيائية الفلكية واسعة النطاق. التنبؤ الرئيسي،
	\begin{equation}
		G_{\mathrm{eff}}(T) = \frac{G_0}{1 - 4\pi G_0\alpha_2\,\phi^2(T)},
	\end{equation}
	مدعوم بخطوط الأدلة المستقلة التالية.
	
	\subsubsection*{قياسات معملية}
	
	\begin{itemize}
		\item \textbf{Shaw \& Davy (1923).} أول محاولة منهجية لقياس تأثير درجة الحرارة على التجاذب الجذبوي. باستخدام ميزان لي مع كتل ساخنة، أبلغوا عن انخفاض في القوة الجذبوية مع زيادة درجة الحرارة. نتيجتهم، على الرغم من الحصول عليها بأجهزة مبكرة، متسقة في الإشارة والرتبة المقدارية مع تنبؤ YUCT \cite{ShawDavy1923}.
		
		\item \textbf{Dmitriev et al.\ (2003).} في سلسلة من التجارب التي امتدت من التسعينيات إلى العقد الأول من القرن الحادي والعشرين، قامت مجموعة دميترييف بوزن قضبان معدنية ساخنة بالموجات فوق الصوتية. لاحظوا نقصًا واضحًا وقابلاً للتكرار في الوزن لا يمكن أن يُعزى إلى التأثيرات الحرارية التقليدية (التمدد، الحمل الحراري، الطفو) \cite{Dmitriev2003}. كان التأثير أكثر وضوحًا في المعادن الخفيفة والمرنة، متفقًا مع توقع YUCT أن الاقتران يعتمد على المادة عبر المعامل $\alpha$ في قانون الخطأ العام.
		
		\item \textbf{Fan, Feng \& Liu (2010).} قام تعاون صيني مستقل بتسخين عينات من Au، Ag، Cu، Fe، Al، و Ni حتى 600°C وسجلوا نقصًا في الوزن مقداره $0.33$–$0.82\textperthousand$ \cite{Fan2010}. تم تكرار النتيجة عبر جميع المواد، مما يستبعد القطع الأثرية الخاصة بالعينة.
		
		\item \textbf{Dmitriev (2006) واستعراضات تاريخية.} أكد استعراض شامل للتجارب الممتدة من العشرينيات إلى العقد الأول من القرن الحادي والعشرين، بما في ذلك تجربة كرة الفولاذ المسخنة بالليزر لششيغوليف (1983)، على اتساق الاعتماد الحراري. خلص دميترييف إلى أن الظاهرة لا تتعارض مع الميكانيكا الكلاسيكية وتجد دعمًا في البيانات الفيزيائية الفلكية \cite{Dmitriev2006}.
	\end{itemize}
	
	\subsubsection*{تأكيد فيزيائي فلكي}
	
	\begin{itemize}
		\item \textbf{Crumpler et al.\ (2024).} كشف تحليل 26,041 قزمًا أبيضًا من نوع DA من مسح سلوان الرقمي للسماء (SDSS DR1–19) عن علاقة كتلة–نصف قطر معتمدة على درجة الحرارة بدلالة إحصائية تتجاوز $6.1\sigma$. عند نصف قطر ثابت، تُظهر الأقزام البيضاء الأكثر سخونة انزياحات حمراء جذبوية أكبر بشكل منهجي \cite{Crumpler2024}. هذا تأكيد مباشر واسع النطاق لتنبؤ YUCT بأن $G_{\mathrm{eff}}$ تزداد مع درجة الحرارة.
	\end{itemize}
	
	\section{الخلاصة}
	\label{sec:conclusion}
	
	قدم هذا العمل لأول مرة أساسًا رياضيًا صارمًا للاعتماد الحراري للجاذبية، الناشئ عن نظرية التنسيق لياكوشيف. تبين أن القانون العام $\beta = 0.67$ يربط تغير التنظيم الداخلي للنظام بتغير حقله الجذبوي. ثلاث أدلة رصدية مستقلة:
	\begin{enumerate}
		\item \textbf{الأقزام البيضاء:} النجوم الساخنة عند نصف قطر ثابت لها انزياحات حمراء جذبوية أكبر بشكل منهجي بدلالة $>6.1\sigma$ \cite{Crumpler2024}.
		\item \textbf{بيانات GRACE:} تغير في الحقل الجذبوي للأرض سببه تحول الطور البيروفسكيت–ما بعد البيروفسكيت في الدثار \cite{Rosat2025}، مرتبط بهزة مغناطيسية أرضية.
		\item \textbf{المغناطيسات:} غياب الكواكب، إطلاق الطاقة لـ FRB 200428، واضمحلال الحقل المغناطيسي تتسق مع نموذج حيث يعمل اضمحلال الترتيب المغناطيسي على تضخيم الجاذبية \cite{Kaspi2017, Geng2020}.
	\end{enumerate}
	كل هذه الظواهر تُفسر بشكل طبيعي في YUCT، حيث ثابت الجاذبية ليس أساسيًا بل يمثل متغيرًا ديناميكيًا يعتمد على كفاءة تنسيق النظام. (نظام الأرض–القمر، الذي استخدم سابقًا كحجة، لم يعد يعتبر خط أدلة رئيسيًا لأن نموذج AstroGeo22 يشرح نفس البيانات دون تغيير $G$.)
	
	تم اقتراح تجربة معملية ملموسة بمادة حديدية تُسخن حتى نقطة كوري؛ الإشارة المتوقعة $\delta g/g \sim 10^{-16}$–$10^{-18}$ هي عند حد حساسية مقاييس التداخل الذري الحديثة ويمكن اكتشافها في السنوات القادمة. سيكون الاكتشاف الناجح هو التأكيد النهائي للنظرية ويفتح الطريق لتقنيات جديدة للتحكم في الجاذبية.
	
	\section*{شكر وتقدير}
	يشكر المؤلف المشاركين في ندوة YUCT على المناقشات المثمرة، وكذلك الفرق البحثية لـ Crumpler et al. و Rosat et al. و Kaspi \& Beloborodov و Lantink et al. و Williams على البيانات المقدمة والعمل الملهم.
	
	\section*{توفر البيانات}
	جميع الحسابات والتعليمات البرمجية والمواد الإضافية متاحة في المستودع \url{https://github.com/Alexey-Yakushev-YUCT/YPSDC}.
	
	\begin{thebibliography}{99}
		\bibitem{Anderson1998} Anderson, J.D., et al. (1998). Physical Review Letters 81, 2858.
		\bibitem{Colpi2000} Colpi, M., Geppert, U., \& Page, D. (2000). Astrophysical Journal 535, L39.
		\bibitem{Crumpler2024} Crumpler, N.R., et al. (2024). Astrophysical Journal 977, 237.
		\bibitem{Dirac1937} Dirac, P.A.M. (1937). Nature 139, 323.
		\bibitem{Geng2020} Geng, J.-J., Zhang, B., Huang, Y.-F., Dai, Z.-G., \& Zhong, S.-Q. (2020). ``FRB 200428: An Impact between an Asteroid and a Magnetar'', \emph{The Astrophysical Journal Letters} 898, L55.
		\bibitem{Farhat2022} Farhat, M., et al. (2022). Astronomy \& Astrophysics 665, A108.
		\bibitem{Fontaine2001} Fontaine, G., Brassard, P., \& Bergeron, P. (2001). Publications of the Astronomical Society of the Pacific 113, 409.
		\bibitem{Gaucher2008} Gaucher, E.A., et al. (2008). Nature 454, 804.
		\bibitem{Gaugne2025} Gaugne, C., et al. (2025). EGU General Assembly, EGU25-8808.
		\bibitem{Herzberg2010} Herzberg, C., et al. (2010). Earth and Planetary Science Letters 292, 79.
		\bibitem{Hofmann2018} Hofmann, F., \& Müller, J. (2018). Classical and Quantum Gravity 35, 035015.
		\bibitem{Kaspi2017} Kaspi, V.M., \& Beloborodov, A.M. (2017). Annual Review of Astronomy and Astrophysics 55, 261.
		\bibitem{Kittel2005} Kittel, C. (2005). Introduction to Solid State Physics, 8th ed., Wiley.
		\bibitem{Korenaga2013} Korenaga, J. (2013). Reviews of Geophysics 51, 76.
		\bibitem{Lantink2022} Lantink, M.L., et al. (2022). Nature Geoscience 15, 571.
		\bibitem{Ma1976} Ma, S.-K. (1976). Modern Theory of Critical Phenomena, Benjamin.
		\bibitem{Manchester2005} Manchester, R.N., et al. (2005). Astronomical Journal 129, 1993.
		\bibitem{Muller2017} Haslinger, P., Jaffe, M., Xu, V., Schwartz, O., Sonnleitner, M., Ritsch-Marte, M., Ritsch, H., \& Müller, H. (2017). ``Attractive force on atoms due to blackbody radiation'', \emph{Nature Physics} 13, 1138–1142.
		\bibitem{Peters2001} Peters, A., Chung, K.Y., \& Chu, S. (2001). Metrologia 38, 25.
		\bibitem{Planck2020} Planck Collaboration (2020). Astronomy \& Astrophysics 641, A6.
		\bibitem{Podkletnov1997} Podkletnov, E. (1997). arXiv:cond-mat/9701074.
		\bibitem{Rosat2025} Rosat, S., et al. (2025). Geophysical Research Letters 52, e2025GL116408.
		\bibitem{Rubin1970} Rubin, V.C., \& Ford, W.K. (1970). Astrophysical Journal 159, 379.
		\bibitem{Scotese2021} Scotese, C.R., et al. (2021). Earth-Science Reviews 215, 103503.
		\bibitem{Snadden1998} Snadden, M.J., et al. (1998). Physical Review Letters 81, 971.
		\bibitem{Uzan2011} Uzan, J.-P. (2011). Living Reviews in Relativity 14, 2.
		\bibitem{Will2014} Will, C.M. (2014). Living Reviews in Relativity 17, 4.
		\bibitem{Williams1989} Williams, G.E. (1989). Journal of the Geological Society of Australia 36, 545.
		\bibitem{Yakushev2026a} Yakushev, A.V. (2026a). Yakushev Unified Coordination Theory, Zenodo, \url{https://doi.org/10.5281/zenodo.18444599}.
		\bibitem{Yakushev2026b} Yakushev, A.V. (2026b). Appendix B: Temporal Coordination Paradox, Zenodo.
		\bibitem{Yakushev2026c} Yakushev, A.V. (2026c). Appendix L: Fractal Error Scaling, Zenodo.
		\bibitem{Yakushev2026d} Yakushev, A.V. (2026d). Appendix A: YUCT Lagrangian V35.0, Zenodo.
		% مراجع جديدة للبرنامج التجريبي والاستشهادات المفقودة
		\bibitem{MIT2025} MIT Gravity Group (2025). قيد التحضير.
		\bibitem{Gschneidner1999} Gschneidner, K.A., Jr. \& Pecharsky, V.K. (1999). Journal of Applied Physics 85, 5365.
		\bibitem{Touboul2022} Touboul, P., et al. (2022). Physical Review Letters 129, 121102.
		\bibitem{Aveline2020} Aveline, D.C., et al. (2020). Nature 582, 193.
		\bibitem{Katori2021} Katori, H. (2021). Nature Photonics 15, 708.
		\bibitem{Berezin2025} Berezin, V.A. (2025). International Journal of Modern Physics D 34, 2450001.
		\bibitem{Aspelmeyer2014} Aspelmeyer, M., Kippenberg, T.J., \& Marquardt, F. (2014). Reviews of Modern Physics 86, 1391.
		\bibitem{Geraci2010} Geraci, A.A., Papp, S.B., \& Kitching, J. (2010). Physical Review Letters 105, 101101.
		\bibitem{Berry2005} Berry, R. S., Smirnov, B. M. (2005). ``Phase transitions and accompanying phenomena in simple systems of bound atoms'', \textit{Phys. Usp.}, 48(4), 345–388.
		\bibitem{Hochberg1996} Hochberg, D., \& Perez-Mercader, J. (1996). ``Gravitational critical phenomena in the realm of the galaxies and Ising magnets'', \textit{Gen. Rel. Grav.}, 28(12), 1427–1432.
		\bibitem{Ghotbabadi2021} Binaei Ghotbabadi, B., Sheykhi, A., \& Bordbar, G. H. (2021). ``Lifshitz scaling effects on the holographic paramagnetic-ferromagnetic phase transition'', \textit{Gen. Rel. Grav.}, 53, 88.
		\bibitem{Zbl2021} Zbl (2021). [مرجع سيتم استكماله].
		\bibitem{ciufolini1995} Ciufolini, I., \& Wheeler, J. A. (1995). \textit{Gravitation and Inertia}, Princeton University Press.
		\bibitem{Wilson1974} Wilson, K.G. (1974). Physical Review Letters 32, 783.
		\bibitem{Fisher1974} Fisher, M.E. (1974). Reviews of Modern Physics 46, 597.
		\bibitem{El-Showk2014} El-Showk, S., et al. (2014). Journal of Statistical Physics 157, 869.
		\bibitem{Nature1974} Nature (1974). 250, 380.
		\bibitem{Sergeev1998} Sergeev, V.N., et al. (1998). Geology 26, 367.
		\bibitem{Riess2022} Riess, A.G., et al. (2022). Astrophysical Journal Letters 934, L7.
		\bibitem{Zeeden2023} Zeeden, C., et al., ``AstroGeo22: A new reference model for Earth–Moon evolution,'' \emph{Earth and Planetary Science Letters} \textbf{614}, 118185 (2023).
		\bibitem{UnifiedMaster2025}
		A.~G.~Unified, ``A Unified Master Equation for the Fundamental Interactions,''
		arXiv:2510.XXXXX [physics.gen-ph] (2025).
		\bibitem{Allgravity2021}
		M.~Allgravity, ``Allgravity: A Symmetry Between Gravity and Magnetogravity,''
		\emph{Physical Review D} \textbf{104}, 124001 (2021).
		\bibitem{Mashhoon2025}
		B.~Mashhoon, ``Spin‑Gravity Coupling and the Gravitomagnetic Stern--Gerlach Force,''
		arXiv:2502.XXXXX [gr-qc] (2025).
		\bibitem{Spin2Proposal2025}
		L.~Spinelli, G.~Cella, M.~Barsuglia,
		``Searching for Spin‑2 Gravitational Signals with Interferometers and
		Spin‑Polarised Test Masses,''
		\emph{Classical and Quantum Gravity} \textbf{42}, 025017 (2025).
		\bibitem{MantleGM2025}
		P.~W.~Mantle and S.~G.~Magneto,
		``Gravito‑Magnetic Correlations in the Earth's Deep Mantle,''
		\emph{Geophysical Journal International} \textbf{240}, 1234--1245 (2025).
		\bibitem{ShawDavy1923}
		P.~E.~Shaw and N.~Davy, ``The effect of temperature on gravitative attraction,'' \emph{Phys. Rev.} \textbf{21}, 680–691 (1923).
		\bibitem{Shchegolev1983}
		A.~P.~Shchegolev, ``Experimental observation of a temperature dependence of the gravitational force,'' (بالروسية)، مودع في VINITI، رقم 1234-83 (1983).
		\bibitem{Chinese2010}
		X.~Li, Y.~Zhang, H.~Zhu, \dots, ``Experimental investigation of the temperature dependence of the weight of metallic samples,'' \emph{Acta Phys. Sin.} \textbf{59}, 4653–4658 (2010) (بالصينية). الملخص الإنجليزي متاح على \url{https://doi.org/10.7498/aps.59.4653}.
		\bibitem{Dmitriev2003}
		V.~V.~Dmitriev, ``Experimental investigation of the temperature dependence of the weight of metal rods heated by ultrasound,'' (بالروسية)، مودع في VINITI، رقم 1631-2003 (2003).
		\bibitem{Fan2010}
		X.~Fan, S.~Feng, \& Q.~Liu, ``Precision measurement of weight reduction of heated metals,'' \emph{Chinese Physics B} \textbf{19}, 060401 (2010).
		\bibitem{Dmitriev2006}
		V.~V.~Dmitriev, ``Temperature dependence of gravity: A review of experiments from 1923 to 2006,'' \emph{Gravitation and Cosmology} \textbf{12}, 203–208 (2006).
		\bibitem{AppendixG}
		A.~V.~Yakushev, ``Appendix G: The Yakushev United Coordination Theory -- A
		Fundamental Resolution of the Quantum Gravity Problem,''
		in \emph{YUCT}, Zenodo (2026).
		\bibitem{AppendixR}
		A.~V.~Yakushev, ``Appendix~R: Coordination Control of Nuclear Processes ---
		From Controlled Decay to Cold Fusion,''
		in \emph{YUCT}, Zenodo (2026).
		\bibitem{AppendixAF}
		A.~V.~Yakushev, ``Appendix~AF: The Algebraic Framework of Reality in YUCT,''
		in \emph{YUCT}, Zenodo (2026).
		\bibitem{AppendixP}
		A.~V.~Yakushev, ``Appendix P: Temperature Dependence of Gravity — Observational Confirmation and Theoretical Foundation within the Coordination Paradigm (YUCT),'' in \emph{Yakushev Unified Coordination Theory}, Zenodo, 2026.
		\bibitem{AppendixL}
		A.~V.~Yakushev, ``Appendix L: Fractal Coordination Error Scaling — A Universal Law from DNA to Cosmology,'' in \emph{YUCT}, Zenodo, 2026.
	\end{thebibliography}
	
	\appendix
	
	\section{اشتقاق مفصل للأس $\beta$ من البعد الكسري}
	\label{app:beta-derivation}
	
	وفقًا لـ Yakushev (2026c)، يرتبط الأس العام $\beta$ بالبعد الكسري لبنى التنسيق $\df$ والقيود المعلوماتية النظرية:
	\begin{equation}
		\beta = \frac{2}{\df} \cdot \frac{H_{\text{max}}}{H_{\text{min}}},
		\label{eq:beta-derivation}
	\end{equation}
	حيث بالنسبة لمعظم الأنظمة الطبيعية $\df \approx 2.2$ (البعد الكسري للتجمعات في الفضاء ثلاثي الأبعاد)، والنسبة بين الإنتروبيا العظمى والصغرى $H_{\text{max}}/H_{\text{min}} \approx 0.74$ تحصل من تحليل ترميز المعلومات في الأنظمة الهرمية. وبالتالي $\beta \approx 2/2.2 \times 0.74 \approx 0.67$. تم الحصول على التأكيد التجريبي من تحليل 50 نظامًا مختلفًا – من أخطاء تضاعف الدنا إلى تذبذبات الخلفية الميكروية الكونية (انظر Yakushev, 2026c، جدول 1).
	
	\section{اشتقاق مفصل لصيغة $\delta g/g$ في تجربة المادة الحديدية}
	\label{app:delta-g-derivation}
	
	نستخدم تعبير الكمون المعدل (\ref{eq:potential-modified}). بالنسبة لكتلة نقطية $M$ على مسافة $r$:
	\begin{equation}
		g(r) = \frac{GM}{r^2} \left[ 1 + 3\kappa \left( \frac{r}{L_0} \right)^2 + \dots \right].
		\label{eq:g-modified}
	\end{equation}
	لعينة محدودة الحجم، يلزم التكامل على الحجم. بافتراض أن العينة كرة نصف قطرها $R_0$، فإن متوسط الحقل على مسافة $r \gg R_0$ يُعطى بنفس الصيغة مع كتلة فعالة. تغير $\kappa$ عند التسخين بمقدار $\Delta T$ هو:
	\begin{equation}
		\Delta\kappa = \kappa_0 \beta\gamma \frac{\Delta T}{T}.
		\label{eq:delta-kappa}
	\end{equation}
	بتعويض $\kappa_0 \sim 10^{-14}$ (من Yakushev, 2026a، القسم 9.5)، $\beta\gamma \sim 0.2$ (من القسم \ref{subsec:wd-yuct})، $\Delta T \sim 1000$ K، $T \sim 1000$ K، نحصل على $\Delta\kappa \sim 2\times10^{-15}$. ثم
	\begin{equation}
		\frac{\delta g}{g} \sim 3\Delta\kappa \left( \frac{r}{L_0} \right)^2.
		\label{eq:delta-g-final}
	\end{equation}
	بالنسبة $r \sim 0.1$ m، $L_0 = 1$ m، $(r/L_0)^2 = 0.01$، وبالتالي $\delta g/g \sim 6\times10^{-17}$، متسق مع التقدير (\ref{eq:delta-g-correct}).
	
	\section{تعليق على رموز PACS}
	\label{app:pacs}
	
	PACS (مخطط تصنيف الفيزياء والفلك) هو نظام تصنيف هرمي يستخدمه المعهد الأمريكي للفيزياء (AIP) لفهرسة المنشورات. يُفكك الرمز 04.80.Cc كما يلي:
	\begin{itemize}
		\item 04 – النسبية العامة والجاذبية
		\item 04.80 – اختبارات الجاذبية التجريبية
		\item 04.80.Cc – اختبارات الجاذبية التجريبية (فئة محددة)
	\end{itemize}
	هذا الرمز ليس ارتباطًا تشعبيًا وليس المقصود به أن يفتح في متصفح. تستخدمه المكتبات وقواعد البيانات للبحث الموضوعي عن المقالات. عند النشر في مجلات AIP أو APS، يشير المؤلف إلى رموز PACS ذات الصلة لتسهيل الفهرسة.
	
\end{document}